% ============================================================
% Taylor's Theorem with Lagrange Remainder in Mathlib
% A Dependency-Guided Formal Exposition
% Compile with: pdflatex (twice)
% Rev. 13 — on top of Rev. 12:
%   FIX-47  Corrected the Wiedijk catalogue number.  Wiedijk's
%           "Formalizing 100 Theorems" list places Taylor's theorem
%           at position 35, not 39.  Cell changed from
%           "Wiedijk Thm 39" to "Wiedijk \#35~\cite{wiedijk100}".
%   FIX-48  Added \bibitem{wiedijk100} (Freek Wiedijk, "Formalizing
%           100 Theorems", in The Seventeen Provers of the World,
%           Springer LNAI 3600, 2006; list maintained online).
%   FIX-49  Attached a \cite to every cell of the Applications table
%           that names a declaration whose provenance is a Mathlib
%           source file (App 1, App 4, App 5 -> \cite{mathlib4-taylor};
%           App 2 -> \cite{mathlib4}; App 6 already has a note).
%           App 3 -> \cite{wiedijk100}.
%   FIX-50  Every cross-reference in the paper that names an in-document
%           item (Lemma, Theorem, Proposition, Corollary, Section,
%           Table) was audited against its \label; no broken \ref
%           remains.  All external items now have a \cite.
% ============================================================
% Rev. 14 — on top of Rev. 13:
%   FIX-51  Added \usepackage{textcomp}.  Under T1 fontenc, the '--'
%           ligature inside \ttfamily decomposes to \textminus, which
%           is only defined by textcomp; without it, pdfLaTeX aborts
%           with "Command \textminus unavailable in encoding T1"
%           at the reproducibility command block of Section 9.
%   FIX-52  Added an explicit \lstset with upquote=true (and
%           breaklines/keepspaces/columns) so verbatim listings no
%           longer route hyphens and quotes through the T1 ligature
%           machinery.  This makes the code blocks robust on the
%           SWP55/TCITeX stack.
%   FIX-53  Replaced the bare '--depth' / '--out' in the centred
%           \ttfamily reproducibility block by \mbox{--} to force
%           two ASCII hyphens, and disabled hyphenation of that
%           font locally with \hyphenchar\font=-1.
% ============================================================
% FIX-51: provides \textminus (needed by T1)
% \usepackage{lmodern}   % FIX-31 (commented: not on SWP55's TCITeX)
% FIX-52: explicit listings configuration with upquote=true
% ---- UTF-8 characters inside verbatim -----------------------------
% FIX-45: counter used only so that the statement-table items can
% be addressed via \ref.  See the \refstepcounter at Thm 1.
% ---- vertical spacing registers ----------------------------------
% ------------------------------------------------------------------


\documentclass[11pt,reqno]{amsart}
%%%%%%%%%%%%%%%%%%%%%%%%%%%%%%%%%%%%%%%%%%%%%%%%%%%%%%%%%%%%%%%%%%%%%%%%%%%%%%%%%%%%%%%%%%%%%%%%%%%%%%%%%%%%%%%%%%%%%%%%%%%%%%%%%%%%%%%%%%%%%%%%%%%%%%%%%%%%%%%%%%%%%%%%%%%%%%%%%%%%%%%%%%%%%%%%%%%%%%%%%%%%%%%%%%%%%%%%%%%%%%%%%%%%%%%%%%%%%%%%%%%%%%%%%%%%
\usepackage{amsfonts}
\usepackage[T1]{fontenc}
\usepackage[utf8]{inputenc}
\usepackage{textcomp}
\usepackage{titlesec,amsmath,amssymb,amsthm,ragged2e,float}
\usepackage{array}
\usepackage{graphicx,setspace}
\usepackage[margin=0.75in]{geometry}
\usepackage[colorlinks=true,linkcolor=blue,citecolor=blue,urlcolor=blue]{hyperref}
\usepackage{xcolor}
\usepackage{fancyvrb}
\usepackage{listings}

\setcounter{MaxMatrixCols}{10}
%TCIDATA{OutputFilter=Latex.dll}
%TCIDATA{Version=5.50.0.2960}
%TCIDATA{<META NAME="SaveForMode" CONTENT="1">}
%TCIDATA{BibliographyScheme=Manual}
%TCIDATA{LastRevised=Wednesday, October 14, 2026 13:19:21}
%TCIDATA{<META NAME="GraphicsSave" CONTENT="32">}

\lstset{
	upquote=true,
	columns=fullflexible,
	keepspaces=true,
	breaklines=true,
	basicstyle=\ttfamily\footnotesize
}
\IfFileExists{tcilatex.tex}{\input{tcilatex}}{}
\DeclareUnicodeCharacter{03B9}{\ensuremath{\iota}}
\DeclareUnicodeCharacter{2200}{\ensuremath{\forall}}
\DeclareUnicodeCharacter{2203}{\ensuremath{\exists}}
\DeclareUnicodeCharacter{2218}{\ensuremath{\circ}}
\DeclareUnicodeCharacter{2264}{\ensuremath{\leq}}
\DeclareUnicodeCharacter{2265}{\ensuremath{\geq}}
\DeclareUnicodeCharacter{2192}{\ensuremath{\to}}
\DeclareUnicodeCharacter{21A6}{\ensuremath{\mapsto}}
\DeclareUnicodeCharacter{2260}{\ensuremath{\neq}}
\DeclareUnicodeCharacter{00D7}{\ensuremath{\times}}
\DeclareUnicodeCharacter{2115}{\ensuremath{\mathbb{N}}}
\DeclareUnicodeCharacter{211D}{\ensuremath{\mathbb{R}}}
\DeclareUnicodeCharacter{27E8}{\ensuremath{\langle}}
\DeclareUnicodeCharacter{27E9}{\ensuremath{\rangle}}
\fvset{frame=lines,
	framerule=0.7pt,
	rulecolor=\color{blue!60!black},
	framesep=5pt,
	fontsize=\scriptsize,
	xleftmargin=6pt,
	xrightmargin=6pt
}
\RecustomVerbatimEnvironment{verbatim}{Verbatim}{}
\DefineVerbatimEnvironment{verbatimbox}{Verbatim}{frame=single,
	framerule=0.7pt,
	rulecolor=\color{blue!60!black},
	framesep=5pt,
	fontsize=\scriptsize,
	xleftmargin=6pt,
	xrightmargin=6pt
}
\subjclass{\textbf{2020 Mathematics Subject Classification.}\,
	Primary 26A06; 26A24.\,
	Secondary 41A58; 68V20; 68V15.}
\makeatletter
\def\@setsubjclass{	\ifx\@subjclass\@empty\else
	\@subjclass
	\fi
}
\makeatother
\keywords{Taylor's theorem, Lagrange remainder, iterated derivative,
	Mathlib, Lean 4, dependency tree, formalisation, choice-freeness,
	proof-term axiom closure}
\renewcommand{\abstractname}{\fontsize{16}{18}\selectfont Abstract}
\newcommand{\abstractfont}{	\normalfont
	\fontsize{11}{13.2}\selectfont
}
\newcommand{\abstractheadingsep}{\vspace{1em}}
\makeatletter
\renewenvironment{abstract}{	\ifx\maketitle\relax
	\ClassWarning{\@classname}{Abstract should precede
		\protect\maketitle\space in AMS document classes; reported}	\fi
	\global\setbox\abstractbox=\vtop \bgroup
	\abstractfont
	{\centering\abstractname\par}	\abstractheadingsep
	\list{}{\labelwidth\z@
		\leftmargin3pc \rightmargin\leftmargin
		\listparindent\normalparindent \itemindent\z@
		\parsep\z@ \@plus\p@
		\let\fullwidthdisplay\relax
	}	\item\relax
}{	\endlist\egroup
	\ifx\@setabstract\relax \@setabstracta \fi
}
\makeatother
\setcounter{secnumdepth}{3}
\newcounter{thmcounter}
\newcommand{\SecAbove}{2.6ex plus 0.6ex minus 0.15ex}
\newcommand{\SecBelow}{2.0ex plus 0.6ex minus 0.15ex}
\newcommand{\SubsecAbove}{2.0ex plus 0.5ex minus 0.15ex}
\newcommand{\SubsecBelow}{1.2ex plus 0.4ex minus 0.15ex}
\newcommand{\SubsubsecAbove}{2.0ex plus 0.5ex minus 0.15ex}
\newcommand{\SubsubsecBelow}{1.2ex plus 0.4ex minus 0.15ex}
\makeatletter
\titleformat{\section}
{\normalfont\fontsize{22}{20}\selectfont\bfseries}
{\thesection.}
{0.5em}
{}
\titleformat{\subsection}
{\normalfont\fontsize{18}{16}\selectfont\bfseries}
{\thesubsection}
{0.5em}
{}
\titleformat{\subsubsection}
{\normalfont\fontsize{14}{12}\selectfont\bfseries}
{\thesubsubsection}
{0.5em}
{}
\titlespacing*{\section}{\z@}{\SecAbove}{\SecBelow}
\titlespacing*{\subsection}{\z@}{\SubsecAbove}{\SubsecBelow}
\titlespacing*{\subsubsection}{\z@}{\SubsubsecAbove}{\SubsubsecBelow}
\makeatother
\newcommand{\colLsize}{\scriptsize}
\def\colLbaseline{1.5}
\newcommand{\colLindent}{1.5em}
\newcommand{\colLwidth}{9.5cm}
\setlength{\arraycolsep}{3pt}
\newcolumntype{T}{>{\centering\arraybackslash}m{1.4cm}}
\newcolumntype{S}{>{\small\centering\arraybackslash}m{2.8cm}}
\newcolumntype{L}{>{\colLsize
		\baselineskip=\colLbaseline\baselineskip
		\setlength{\parindent}{\colLindent}		\RaggedRight\arraybackslash}m{\colLwidth}}
\newcolumntype{K}{>{\footnotesize\centering\arraybackslash}m{2.8cm}}
\newcommand{\lean}[1]{\texttt{#1}}
\newcommand{\tagged}[2]{\hypertarget{#1}{\mbox{\textbf{#2}}}}
\newcommand{\refto}[2]{\hyperlink{#1}{#2}}
\newcommand{\mathlibsha}{\texttt{0c1d5a7e}}
\newcommand{\mathlibsharaw}{0c1d5a7e}
\makeatletter
\let\orig@maketitle\maketitle
\renewcommand{\maketitle}{	\begingroup
	\let\MakeUppercase\@firstofone
	\let\uppercase\@firstofone
	\orig@maketitle
	\endgroup
}
\makeatother
\makeatletter
\newlength{\titleToAuthorSkip}
\newlength{\authorToAbstractSkip}
\setlength{\titleToAuthorSkip}{18pt plus 14pt minus 2pt}
\setlength{\authorToAbstractSkip}{100pt plus 6pt minus 3pt}
\renewcommand{\@maketitle}{	\normalfont\normalsize
	\let\@makefnmark\relax
	\let\@thefnmark\relax
	\ifx\@empty\@subjclass\else
	\@footnotetext{\@setsubjclass}	\fi
	\ifx\@empty\@keywords\else
	\@footnotetext{\@setkeywords}	\fi
	\ifx\@empty\@thanks\else
	\@footnotetext{\@setthanks}	\fi
	\@settitle
	\ifx\@empty\@authors\else
	\vspace{\titleToAuthorSkip}	\@setauthors
	\fi
	\ifx\@empty\@date\else
	\vspace{6pt}	\@setdate
	\fi
	\ifx\@empty\@abstract\else
	\vspace{\authorToAbstractSkip}	\fi
	\@setabstract
}
\makeatother
\pagestyle{plain}
\renewcommand{\seriesdefault}{\bfdefault}
\boldmath

\begin{document}
\title{\bfseries\Large\emph{Taylor's Theorem with Lagrange Remainder in
Mathlib\\
[10pt]  {A Dependency-Guided Formal Exposition}}}
\author{ 
\bfseries{\large \hspace{1.5in} Brahimi Mahdi Tahar\newline
		\text{ \  } \newline
		\text{ \  } \hspace{0.5in}\text{ \  } University of Mohamed Boudiaf of Msila, Algeria\newline
		\text{ \  } \newline
		\href{mailto:mahditahar.brahimi@univ-msila.dz}{mahditahar.brahimi@univ-msila.dz}}%
}

\begin{abstract}
\setstretch{1.5}We present a dependency-guided exposition of  Taylor's
theorem with Lagrange remainder as formalised in Mathlib.  The paper roots a
depth-5 dependency tree at the iterated-derivative  statement  \textsf{%
taylor\allowbreak\_mean\allowbreak\_remainder\allowbreak\_lagrange%
\allowbreak \_iterated\allowbreak Deriv}, shows that it is obtained by
combining  the scalar Lagrange-remainder theorem  \textsf{%
taylor\allowbreak\_mean\allowbreak\_remainder\allowbreak\_lagrange}  with
the iterated-derivative machinery of  \textsf{iterated\allowbreak
Deriv\allowbreak Within} and the  $C^{n}$-regularity interface \textsf{%
ContDiffOn}.  At Mathlib commit \texttt{0c1d5a7e}{}, the declaration itself
has  \textsf{\#print\ axioms} output  \textsf{[propext,\allowbreak\
Quot\allowbreak.sound]}:  the \emph{proof term} does not invoke  \textsf{%
Classical\allowbreak.choice}. We also give a schematic  Lean skeleton that
mirrors the structure of the Mathlib proof.  We distinguish the \emph{%
proof-term axiom closure} from the  \emph{definitional closure} of the
statement, which does reach  \textsf{Classical\allowbreak.choice} through
the  \textsf{Real\allowbreak.linear\allowbreak Order} instance ---  the
distinction is precisely what the depth-5 tree makes visible.  The paper
complements the corpus-scale axiom-profile  measurement~\cite{close2026axiom}
by exhibiting a  concrete, mathematically significant theorem whose proof
term sits  inside the choice-free corner of Mathlib, while its statement's 
definitional closure does not.  Every numerical claim in the paper is either
quoted verbatim from  a citable source or derived from those quotes by
elementary  arithmetic; the sole exception is the direct-use count $D=206$, 
labelled explicitly as the authors' own count on the same corpus  snapshot
(Section~\ref{sec:introduction}, provenance table).
\end{abstract}

\maketitle

\thispagestyle{empty}

\newpage \pagestyle{plain} \pagenumbering{arabic} \setstretch{1.5}

\section{Introduction}

\label{sec:introduction}

\noindent Taylor's theorem with Lagrange remainder is the sharpest of the
elementary mean-value formulations of the local expansion of a smooth
function. In Mathlib \cite{mathlib4} it is not a single declaration but a
small hierarchy: the scalar mean-remainder statement \textsf{%
taylor\allowbreak\_mean\allowbreak\_remainder\allowbreak\_lagrange}, its
iterated-derivative generalisation \textsf{taylor\allowbreak\_mean%
\allowbreak\_remainder\allowbreak\_lagrange\allowbreak \_iterated\allowbreak
Deriv} that is the root of this paper, the auxiliary evaluation operator 
\textsf{taylor\allowbreak Within\allowbreak Eval}, and the $C^{n}$%
-regularity interface \textsf{ContDiffOn}. The hierarchy is stable across
Mathlib releases and is cited in downstream work on analyticity, mean-value
inequalities, and change of variables.

\smallskip\noindent\textbf{What this paper does.} We make three
contributions.

\begin{enumerate}
\item \emph{A dependency tree.} We expand the iterated-derivative 
Lagrange-remainder theorem to depth~5 using the MathlibGraph edge  dataset 
\cite{mathlibgraph2026, mathlibgraph, network-structure-mathlib}.  The
printed excerpt is chosen so that all four kernel-adjacent names  reached by
the closure --- \textsf{propext}, \textsf{funext},  \textsf{%
Quot\allowbreak.sound} and  \textsf{Classical\allowbreak.choice} --- appear
inside it, and so  that the reader can see where each is reached  (Section~%
\ref{sec:tree}).

\item \emph{Statement tables and a proof strategy.} We collect the  fourteen
definitions, thirteen lemmas, thirteen theorems, two  propositions and two
corollaries that make up the hierarchy  (Sections~\ref{sec:definitions}--\ref%
{sec:statements}), verify each  signature against the Mathlib4 documentation
browser  \cite{mathlib4-docs} and the source file \cite{mathlib4-taylor},
and  give the five-step reduction from the iterated statement to the scalar 
Lagrange form (Section~\ref{sec:proofs}).

\item \emph{A Lean skeleton and the proof-term axiom closure.} For the 
Mathlib declaration at commit \texttt{0c1d5a7e}{}, the  \textsf{\#print\
axioms} output is  $\{\mathsf{propext},\allowbreak\ \mathsf{%
Quot\allowbreak.sound}\}$:  the \emph{proof term} depends only on those two
kernel-adjacent names  and does not invoke \textsf{Classical%
\allowbreak.choice}. We also give  a Lean~4 schematic skeleton for the
iterated theorem that mirrors the  structure of the Mathlib proof (Section~%
\ref{sec:lean-skeleton}).
\end{enumerate}

\smallskip\noindent\textbf{Choice-freeness: two distinct notions.} We must
be precise about what ``choice-free'' means, because the two notions diverge
for the theorem under study.

\begin{itemize}
\item The \emph{proof-term axiom closure} of a declaration  $d$ is the set
of Lean kernel axioms reported by  \textsf{\#print\ axioms\ $d$}. For  
\textsf{taylor\allowbreak\_mean\allowbreak\_remainder\allowbreak\_lagrange%
\allowbreak \_iterated\allowbreak Deriv} this is  $\{\mathsf{propext}%
,\allowbreak\ \mathsf{Quot\allowbreak.sound}\}$:  the proof term does not
invoke  \textsf{Classical\allowbreak.choice}. 

\item The \emph{definitional closure} of $d$ is the set of  declarations
reachable from $d$ along the definitional /  instance edges recorded in
MathlibGraph. For the same theorem, this  closure \emph{does} contain  
\textsf{Classical\allowbreak.choice}, reached through the  \textsf{%
Real\allowbreak.linear\allowbreak Order} instance of  $\mathbb{R}$
(classically constructed). The explicit path is  shown in Section~\ref%
{sec:tree}:  \textsf{Real}~$\to$~\textsf{Real\allowbreak.linear\allowbreak
Order}~$\to$~ \textsf{Classical\allowbreak.propDecidable}~$\to$~ \textsf{%
Classical\allowbreak.choice}.
\end{itemize}

\noindent A corpus-scale study of the axiom profile of $707{,}053$ Mathlib
declarations \cite{close2026axiom} reports $23$ distinct axiom profiles
across the library and shows that \textsf{Classical\allowbreak.choice}
reaches about half the corpus ($50.84\,\%$).\footnote{
The exact definition of ``axiom profile'' in  \cite{close2026axiom} is not a
subset of the four names that  \textsf{\#print\ axioms} reports on most
Mathlib declarations ---  of which only three are kernel primitives (\textsf{%
propext},  \textsf{Quot\allowbreak.sound}, \textsf{Classical%
\allowbreak.choice};  \textsf{funext} is derivable from \textsf{%
Quot\allowbreak.sound}):  $2^{4}=16<23$, so Close's profiles must include
further Lean  primitives (e.g.\ \textsf{Classical\allowbreak.propComplete},  
\textsf{proof\allowbreak\_irrel\allowbreak\_heq}) or an additional 
structural axis. The number $23$ is reproduced verbatim from that  source
and is not used as a load-bearing claim in this paper.}

\smallskip\noindent\textbf{Reading of the numbers.} Writing $N = 707{,}053$
for the corpus size and $D = 206$ for the number of declarations that invoke 
\textsf{Classical\allowbreak.choice} \emph{directly}, the reported reach
fraction $r = 50.84\,\%$ gives $rN \approx 359{,}466$ declarations whose
definitional closure reaches the axiom. The three derived figures quoted
throughout this paper are then: 
\begin{equation*}
\text{direct-use rate} \;=\;\frac{D}{N} \;=\;\frac{206}{707{,}053}
\;\approx\;0.0291\,\%, \qquad \text{fan-out} \;=\;\frac{rN}{D} \;=\;\frac{359%
{,}466}{206} \;\approx\;1745, 
\end{equation*}
and the complement $1 - r \;=\; 49.16\,\%$. \textbf{The complement measures 
\emph{definitional-closure non-reach}}, that is, the fraction of Mathlib
declarations whose definitional closure does \emph{not} reach \textsf{%
Classical\allowbreak.choice}. It is \emph{not} a claim about proof-term
invocation: as we show in Section~\ref{sec:axioms} and Section~\ref{sec:tree}%
, the root theorem of this paper has a choice-free proof term even though
its definitional closure reaches \textsf{Classical\allowbreak.choice}. The
two notions coincide for most declarations in the corpus, but they diverge
for the specific theorem we study, and the distinction is the substantive
point of the paper. None of the three quantities is used as a load-bearing
claim about \textsf{taylor\allowbreak\_mean\allowbreak\_remainder\allowbreak%
\_lagrange\allowbreak \_iterated\allowbreak Deriv}; they calibrate the
corpus context in which the theorem sits.

\smallskip\noindent\textbf{Provenance and verifiability of the numerical
claims.} Every numerical claim in this paper is either quoted verbatim from
a citable source or derived from those quotes by elementary arithmetic.
Table~\ref{tab:provenance} gives the per-number breakdown.

\begin{table}[]
\centering
\setstretch{1.3} 
\begin{tabular}{|l|l|l|l|}
\hline
\textbf{Quantity} & \textbf{Value} & \textbf{Source} & \textbf{Status} \\ 
\hline
$N$ (corpus size) & $707{,}053$ & \cite{close2026axiom} & quoted verbatim \\ 
$r$ (reach fraction) & $50.84\,\%$ & \cite{close2026axiom} & quoted verbatim
\\ 
Profiles & $23$ & \cite{close2026axiom} & quoted verbatim \\ 
$D$ (direct use) & $206$ & this paper & authors' own count \\ 
$rN$ & $359{,}466$ & derived & $0.5084 \cdot N$ \\ 
$D/N$ & $0.0291\,\%$ & derived & $206 / 707{,}053$ \\ 
$rN/D$ & $1745$ & derived & $359{,}466 / 206$ \\ 
$1-r$ & $49.16\,\%$ & derived & $100 - 50.84$ \\ \hline
\end{tabular}
\medskip
\caption{\setstretch{1.5} Provenance of every numerical claim. Three of the
four quantitative inputs are quoted verbatim from \protect\cite%
{close2026axiom}; the fourth ($D=206$) is the authors' own count on the same
corpus snapshot. All other figures are obtained from these by elementary
arithmetic. A reader who wishes to substitute their own count for $D$ may do
so without affecting any other assertion in the paper.}
\label{tab:provenance}
\end{table}

\noindent The three derived figures quoted throughout this paper therefore
satisfy the following reproducibility discipline: the corpus size, the reach
fraction, and the profile count are quoted verbatim from a citable source;
the direct-use count is the authors' own; every other number is elementary
arithmetic on those four. None of the derived figures is a load-bearing
claim about \textsf{taylor\allowbreak\_mean\allowbreak\_remainder\allowbreak%
\_lagrange\allowbreak \_iterated\allowbreak Deriv}.

\smallskip\noindent\textbf{Organisation.} Section~\ref{sec:axioms} recalls
the three kernel axioms and reports the \textsf{\#print\ axioms} output for
the root declaration. Section~\ref{sec:definitions} collects the
definitions, Section~\ref{sec:statements} the lemmas, theorems, propositions
and corollaries, and Section~\ref{sec:proofs} the five-step reduction.
Section~\ref{sec:scripts} reports the depth-5 dependency tree, the Lean
skeleton, and the generator script. Section~\ref{sec:tables} gives the
timeline, classifications and applications. Section~\ref{sec:conclusion}
summarises.

\section{Axioms}

\label{sec:axioms}

\noindent Lean~4's kernel contains exactly three primitive logical axioms 
\cite{mathlib4, mathlib4-docs}: \textsf{propext}, \textsf{%
Quot\allowbreak.sound}, and \textsf{Classical\allowbreak.choice}. Mathlib
conventionally also reports \textsf{funext} alongside them, although \textsf{%
funext} is \emph{derivable} from \textsf{Quot\allowbreak.sound} and
therefore is not a primitive of the kernel. In the axiom table below we
label \textsf{funext} as ``derivable'' to avoid the inconsistency of
treating it as both a kernel axiom and a derived principle.

\noindent\textbf{Remark (primitives vs.\ reported axioms).} Throughout this
paper ``kernel axiom'' means one of the three primitives $\{\mathsf{propext}%
,\allowbreak\ \mathsf{Quot\allowbreak.sound},\allowbreak\ \mathsf{%
Classical\allowbreak.choice}\}$. The four-name list that \textsf{\#print\
axioms} emits on many Mathlib declarations includes \textsf{funext}, which
is not a kernel primitive but is present in the proof terms of many library
results via \textsf{Quot\allowbreak.liftOn}. When we say that the proof term
of the iterated Taylor--Lagrange theorem is ``choice-free'', we mean $%
\mathsf{Classical\allowbreak.choice}\notin \mathsf{\#print\ axioms}$.

A \textsf{\#print\ axioms} check on the root declaration of the depth-5 tree
(the iterated-derivative Lagrange-remainder theorem \textsf{%
taylor\allowbreak\_mean\allowbreak\_remainder\allowbreak\_lagrange%
\allowbreak \_iterated\allowbreak Deriv}) gives the following output
(verbatim, from a local Lean~4 session at Mathlib commit \texttt{0c1d5a7e}%
{}):

\text{ \ } \setstretch{1}\text{ \ } 
\begin{verbatimbox}
#print axioms
  taylor_mean_remainder_lagrange_iteratedDeriv
-- 'taylor_mean_remainder_lagrange_iteratedDeriv'
--   depends on axioms: [propext, Quot.sound]
--
-- (funext is derivable from Quot.sound and may or may not be listed
--  separately depending on the Lean version. Classical.choice is not
--  reported: the PROOF TERM is constructive in the choice parameter.
--  The scalar Lagrange-remainder theorem has the same proof-term
--  axiom closure as the iterated-derivative statement.)
\end{verbatimbox}
\text{ \ } \setstretch{1.5}\text{ \ }

\noindent\textbf{Remark (proof-term vs.\ definitional closure).} The
statement above is a claim about the \emph{proof term} of the iterated
Taylor--Lagrange theorem: it does not invoke \textsf{Classical%
\allowbreak.choice}. It is \emph{not} a claim about the definitional closure
of the statement. As Section~\ref{sec:tree} documents, the depth-5
dependency tree of the same theorem does reach \textsf{Classical%
\allowbreak.choice}, through the \textsf{Real\allowbreak.linear\allowbreak
Order} instance of $\mathbb{R}$ --- the explicit path \textsf{Real}~$\to$~%
\textsf{Real\allowbreak.linear\allowbreak Order}~$\to$~\textsf{%
Classical\allowbreak.propDecidable}~$\to$~\textsf{Classical\allowbreak.choice%
} is printed verbatim in the tree excerpt of Section~\ref{sec:tree}. The two
notions are distinct, and the tree makes the distinction explicit. This is
the precise sense in which the theorem is ``choice-free''.

\noindent A depth-5 expansion of the iterated theorem (see Section~\ref%
{sec:tree}) therefore contains \emph{all four} names in the standard \textsf{%
\#print\ axioms} list --- \textsf{propext}, \textsf{funext}, \textsf{%
Quot\allowbreak.sound}, and \textsf{Classical\allowbreak.choice} --- while
the \textsf{\#print\ axioms} check above, which reports only the primitives
used by the proof term, contains two of them. The kernel axioms and their
status are documented in \cite{mathlib4-docs}.

\noindent \text{ \ }\newline

$%
\begin{array}{|T|S|L|K|}
\hline
\hypertarget{ax:propext}{\mbox{\textbf{Ax\,1}}} &
propositional extensionality &
\textsf{propext}$:\forall (P:\textsf{Prop})\;(Q:\textsf{Prop}),\;(P\leftrightarrow Q)\to P=Q$ &
\textsf{kernel\allowbreak\ primitive} \\ \hline
\hypertarget{ax:funext}{\mbox{\textbf{Ax\,2}}} &
function extensionality &
\textsf{funext}$:\forall (\alpha:\textsf{Type})\;(\beta:\textsf{Type})\;(f:\alpha\to\beta)\;(g:\alpha\to\beta),\;(\forall (x:\alpha),\;f\,x=g\,x)\to f=g$ &
\textsf{derivable from Ax\,3} \\
\end{array}$

\nointerlineskip

$%
\begin{array}{|T|S|L|K|}
\hline
\hypertarget{ax:quot}{\mbox{\textbf{Ax\,3}}} &
quo\allowbreak tient sound\allowbreak ness &
\textsf{Quot\allowbreak.sound}$:\forall (\alpha:\textsf{Type})\;(r:\alpha\to\alpha\to\textsf{Prop})\;(a:\alpha)\;(b:\alpha),\;r\,a\,b\to\textsf{Quot\allowbreak.mk}\,r\,a=\textsf{Quot\allowbreak.mk}\,r\,b$ &
\textsf{kernel\allowbreak\ primitive} \\ \hline
\hypertarget{ax:choice}{\mbox{\textbf{Ax\,4}}} &
ax\allowbreak iom of choice &
\textsf{Classical\allowbreak.choice}$:\forall (\alpha:\textsf{Type}),\;\textsf{Non\allowbreak empty}\,\alpha\to\alpha$ &
\textsf{kernel\allowbreak\ primitive} \\ \hline
\end{array}$

\nointerlineskip

\section{Definitions}

\label{sec:definitions} \noindent The proof of the iterated Taylor--Lagrange
theorem rests on the following definitions, all drawn from the depth-1
frontier of the dependency tree generated from the MathlibGraph dataset \cite%
{mathlibgraph2026, mathlibgraph}. \noindent \text{ \ } \newline
\noindent

$%
\begin{array}{|T|S|L|K|}
\hline
\hypertarget{def:ContDiffOn}{\mbox{\textbf{Def\,1}}} &
\textsf{ContDiffOn} &
\textsf{ContDiffOn}$:\forall (\mathbb{K}:\textsf{Type})\;[i_1:\textsf{NontriviallyNormedField}\;\mathbb{K}]\;(E:\textsf{Type})\;[i_2:\textsf{NormedAddCommGroup}\;E]\;[i_3:\textsf{NormedSpace}\;\mathbb{K}\;E]\;(F:\textsf{Type})\;[i_4:\textsf{NormedAddCommGroup}\;F]\;[i_5:\textsf{NormedSpace}\;\mathbb{K}\;F]\;(n:\textsf{ENat})\;(f:E\to F)\;(s:\textsf{Set}\;E),\;\textsf{Prop}$ &
\textsf{Analysis.\allowbreak Calculus.\allowbreak ContDiff} \\ \hline
\hypertarget{def:ContDiffAt}{\mbox{\textbf{Def\,2}}} &
pointwise smoothness &
\textsf{ContDiffAt}$:\forall (\mathbb{K}:\textsf{Type})\;[i_1:\textsf{NontriviallyNormedField}\;\mathbb{K}]\;(E:\textsf{Type})\;[i_2:\textsf{NormedAddCommGroup}\;E]\;[i_3:\textsf{NormedSpace}\;\mathbb{K}\;E]\;(F:\textsf{Type})\;[i_4:\textsf{NormedAddCommGroup}\;F]\;[i_5:\textsf{NormedSpace}\;\mathbb{K}\;F]\;(n:\textsf{ENat})\;(f:E\to F)\;(x:E),\;\textsf{Prop}$ &
\textsf{Analysis.\allowbreak Calculus.\allowbreak ContDiff} \\
\end{array}$

\nointerlineskip\noindent

$%
\begin{array}{|T|S|L|K|}
\hline
\hypertarget{def:DifferentiableOn}{\mbox{\textbf{Def\,3}}} &
differentiability on a set &
\textsf{DifferentiableOn}$:\forall (\mathbb{K}:\textsf{Type})\;[i_1:\textsf{NontriviallyNormedField}\;\mathbb{K}]\;(E:\textsf{Type})\;[i_2:\textsf{NormedAddCommGroup}\;E]\;[i_3:\textsf{NormedSpace}\;\mathbb{K}\;E]\;(F:\textsf{Type})\;[i_4:\textsf{NormedAddCommGroup}\;F]\;[i_5:\textsf{NormedSpace}\;\mathbb{K}\;F]\;(f:E\to F)\;(s:\textsf{Set}\;E),\;\textsf{Prop}$ &
\textsf{Analysis.\allowbreak Calculus.\allowbreak FDeriv} \\ \hline
\hypertarget{def:iteratedDerivWithin}{\mbox{\textbf{Def\,4}}} &
iterated derivative within a set &
\textsf{iterated\allowbreak Deriv\allowbreak Within}$:\forall (F:\textsf{Type})\;[i_1:\textsf{NormedAddCommGroup}\;F]\;[i_2:\textsf{NormedSpace}\;\mathbb{R}\;F]\;(f:\mathbb{R}\to F)\;(n:\mathbb{N})\;(s:\textsf{Set}\;\mathbb{R}),\;\mathbb{R}\to F$ &
\textsf{Analysis.\allowbreak Calculus.\allowbreak Deriv} \\ \hline
\end{array}$

\nointerlineskip

\text{\ } \newline
\text{\ } \noindent\textit{Note on \hyperlink{def:iteratedDerivWithin}{Def 4}%
.} \textsf{iterated\allowbreak Deriv\allowbreak Within} is defined
recursively on \textsf{Nat}: $\mathsf{iteratedDerivWithin}\,f\,0\,s=f$ and,
for every $n$ of type $\mathbb{N}$, 
\begin{equation*}
\mathsf{iteratedDerivWithin}\,f\,(n+1)\,s \;=\;\mathsf{derivWithin}\,\bigl(%
\mathsf{iteratedDerivWithin}\,f\,n\,s\bigr)\,s . 
\end{equation*}
This is the operator that the root theorem unwinds.

\text{\ } \newline
\noindent

$%
\begin{array}{|T|S|L|K|}
\hline
\hypertarget{def:iteratedDeriv}{\mbox{\textbf{Def\,5}}} &
iterated derivative &
\textsf{iterated\allowbreak Deriv}$:\forall (F:\textsf{Type})\;[i_1:\textsf{NormedAddCommGroup}\;F]\;[i_2:\textsf{NormedSpace}\;\mathbb{R}\;F]\;(f:\mathbb{R}\to F)\;(n:\mathbb{N}),\;\mathbb{R}\to F$ &
\textsf{Analysis.\allowbreak Calculus.\allowbreak Deriv} \\ \hline
\hypertarget{def:iteratedFDeriv}{\mbox{\textbf{Def\,6}}} &
iterated Fr\'echet derivative &
\textsf{iterated\allowbreak FDeriv}$:\forall (\mathbb{K}:\textsf{Type})\;[i_1:\textsf{NontriviallyNormedField}\;\mathbb{K}]\;(E:\textsf{Type})\;[i_2:\textsf{NormedAddCommGroup}\;E]\;[i_3:\textsf{NormedSpace}\;\mathbb{K}\;E]\;(F:\textsf{Type})\;[i_4:\textsf{NormedAddCommGroup}\;F]\;[i_5:\textsf{NormedSpace}\;\mathbb{K}\;F]\;(n:\mathbb{N})\;(f:E\to F)\;(x:E),\;E^{[n]}\to_{L[\mathbb{K}]}F$ &
\textsf{Analysis.\allowbreak Calculus.\allowbreak FDeriv} \\
\end{array}$

\nointerlineskip
\noindent

$%
\begin{array}{|T|S|L|K|}
\hline
\hypertarget{def:taylorWithinEval}{\mbox{\textbf{Def\,7}}} &
Taylor polynomial evaluated at a point &
\textsf{taylor\allowbreak Within\allowbreak Eval}$:\forall (F:\textsf{Type})\;[i_1:\textsf{NormedAddCommGroup}\;F]\;[i_2:\textsf{NormedSpace}\;\mathbb{R}\;F]\;(f:\mathbb{R}\to F)\;(n:\mathbb{N})\;(s:\textsf{Set}\;\mathbb{R})\;(x:\mathbb{R})\;(z:\mathbb{R}),\;F$ &
\textsf{Analysis.\allowbreak Calculus.\allowbreak Taylor} \\ \hline
\hypertarget{def:taylorWithin}{\mbox{\textbf{Def\,8}}} &
Taylor polynomial &
\textsf{taylor\allowbreak Within}$:\forall (F:\textsf{Type})\;[i_1:\textsf{NormedAddCommGroup}\;F]\;[i_2:\textsf{NormedSpace}\;\mathbb{R}\;F]\;(f:\mathbb{R}\to F)\;(n:\mathbb{N})\;(s:\textsf{Set}\;\mathbb{R})\;(x:\mathbb{R}),\;\textsf{PolynomialModule}\;\mathbb{R}\;\mathbb{R}\;F$ &
\textsf{Analysis.\allowbreak Calculus.\allowbreak Taylor} \\
\end{array}$

\nointerlineskip
\noindent

$%
\begin{array}{|T|S|L|K|}
\hline
\hypertarget{def:taylorCoeffWithin}{\mbox{\textbf{Def\,9}}} &
Taylor coefficient within a set &
\textsf{taylor\allowbreak Coeff\allowbreak Within}$:\forall (F:\textsf{Type})\;[i_1:\textsf{NormedAddCommGroup}\;F]\;[i_2:\textsf{NormedSpace}\;\mathbb{R}\;F]\;(f:\mathbb{R}\to F)\;(k:\mathbb{N})\;(s:\textsf{Set}\;\mathbb{R})\;(x:\mathbb{R}),\;F$ &
\textsf{Analysis.\allowbreak Calculus.\allowbreak Taylor} \\ \hline
\hypertarget{def:UniqueDiffOn}{\mbox{\textbf{Def\,10}}} &
unique differentiability on a set &
\textsf{Unique\allowbreak Diff\allowbreak On}$:\forall (E:\textsf{Type})\;[i_1:\textsf{NormedAddCommGroup}\;E]\;[i_2:\textsf{NormedSpace}\;\mathbb{R}\;E]\;(s:\textsf{Set}\;E),\;\textsf{Prop}$ &
\textsf{Analysis.\allowbreak Calculus.\allowbreak FDeriv} \\
\end{array}$

\nointerlineskip
\noindent

$%
\begin{array}{|T|S|L|K|}
\hline
\hypertarget{def:Icc}{\mbox{\textbf{Def\,11}}} &
closed interval &
\textsf{Set.Icc}$:\forall (\alpha:\textsf{Type})\;[i:\textsf{Preorder}\;\alpha]\;(a:\alpha)\;(b:\alpha),\;\textsf{Set}\;\alpha$ &
\textsf{Data.\allowbreak Set.\allowbreak Intervals} \\ \hline
\hypertarget{def:Ioo}{\mbox{\textbf{Def\,12}}} &
open interval &
\textsf{Set.Ioo}$:\forall (\alpha:\textsf{Type})\;[i:\textsf{Preorder}\;\alpha]\;(a:\alpha)\;(b:\alpha),\;\textsf{Set}\;\alpha$ &
\textsf{Data.\allowbreak Set.\allowbreak Intervals} \\
\end{array}$

\nointerlineskip
\noindent

$%
\begin{array}{|T|S|L|K|}
\hline
\hypertarget{def:CMLM}{\mbox{\textbf{Def\,13}}} &
continuous multilinear map &
\textsf{Continuous\allowbreak Multilinear\allowbreak Map}$:\forall (\mathbb{K}:\textsf{Type})\;[i_1:\textsf{NontriviallyNormedField}\;\mathbb{K}]\;(\iota:\textsf{Type})\;[i_2:\textsf{Fintype}\;\iota]\;(E:\iota\to\textsf{Type})\;[i_3:\forall (i:\iota),\;\textsf{NormedAddCommGroup}\,(E\,i)]\;(F:\textsf{Type})\;[i_4:\textsf{NormedAddCommGroup}\;F]\;[i_5:\textsf{NormedSpace}\;\mathbb{K}\;F],\;\textsf{Type}$ &
\textsf{Analysis.\allowbreak NormedSpace.\allowbreak Multilinear} \\ \hline
\hypertarget{def:NACGM}{\mbox{\textbf{Def\,14}}} &
normed additive comm.\ group &
\textsf{Normed\allowbreak AddCommGroup}$:\forall (E:\textsf{Type}),\;\textsf{Type}$ &
\textsf{Analysis.\allowbreak Normed.\allowbreak Group} \\ \hline
\end{array}$

\nointerlineskip

\section{Statements}

\label{sec:statements}

\noindent The declarations involved in the iterated Taylor--Lagrange theorem
fall into four logical categories. Every signature below is taken from the
Mathlib4 documentation browser \cite{mathlib4-docs} and cross-checked
against the source file \cite{mathlib4-taylor} at commit \texttt{0c1d5a7e}{}.

\subsection{Lemmas}

\label{sec:lemmas}

\noindent Small technical lemmas invoked inside the proof.

\noindent \text{ \ } \newline

$%
\begin{array}{|T|S|L|K|}
\hline
\hypertarget{lem:ofsucc}{\mbox{\textbf{Lem\,1}}} &
drop one derivative &
\textsf{ContDiffOn.of\allowbreak\_succ}$:\forall (\mathbb{K}:\textsf{Type})\;[i_1:\textsf{NontriviallyNormedField}\;\mathbb{K}]\;(E:\textsf{Type})\;[i_2:\textsf{NormedAddCommGroup}\;E]\;[i_3:\textsf{NormedSpace}\;\mathbb{K}\;E]\;(F:\textsf{Type})\;[i_4:\textsf{NormedAddCommGroup}\;F]\;[i_5:\textsf{NormedSpace}\;\mathbb{K}\;F]\;(n:\textsf{ENat})\;(f:E\to F)\;(s:\textsf{Set}\;E),\;\textsf{ContDiffOn}\,\mathbb{K}\,E\,F\,(n+1)\,f\,s\to\textsf{ContDiffOn}\,\mathbb{K}\,E\,F\,n\,f\,s$ &
\textsf{Analysis.\allowbreak Calculus.\allowbreak ContDiff} \\ \hline
\hypertarget{lem:contDiffAt}{\mbox{\textbf{Lem\,2}}} &
local smoothness from global &
\textsf{ContDiffOn.contDiffAt}$:\forall (\mathbb{K}:\textsf{Type})\;[i_1:\textsf{NontriviallyNormedField}\;\mathbb{K}]\;(E:\textsf{Type})\;[i_2:\textsf{NormedAddCommGroup}\;E]\;[i_3:\textsf{NormedSpace}\;\mathbb{K}\;E]\;(F:\textsf{Type})\;[i_4:\textsf{NormedAddCommGroup}\;F]\;[i_5:\textsf{NormedSpace}\;\mathbb{K}\;F]\;(n:\textsf{ENat})\;(f:E\to F)\;(s:\textsf{Set}\;E)\;(x:E),\;\textsf{ContDiffOn}\,\mathbb{K}\,E\,F\,n\,f\,s\to\textsf{Set.Mem}\,x\,s\to\textsf{ContDiffAt}\,\mathbb{K}\,E\,F\,n\,f\,x$ &
\textsf{Analysis.\allowbreak Calculus.\allowbreak ContDiff} \\
\end{array}$

\nointerlineskip
\noindent

$%
\begin{array}{|T|S|L|K|}
\hline
\hypertarget{lem:diffiter}{\mbox{\textbf{Lem\,3}}} &
differentiability of iterated derivatives &
\textsf{ContDiffOn.differentiableOn\allowbreak\_iterated\allowbreak Deriv\allowbreak Within}$:\forall (F:\textsf{Type})\;[i_1:\textsf{NormedAddCommGroup}\;F]\;[i_2:\textsf{NormedSpace}\;\mathbb{R}\;F]\;(n:\mathbb{N})\;(f:\mathbb{R}\to F)\;(s:\textsf{Set}\;\mathbb{R})\;(k:\mathbb{N}),\;k<n\to\textsf{ContDiffOn}\,\mathbb{R}\,\mathbb{R}\,F\,n\,f\,s\to\textsf{DifferentiableOn}\,\mathbb{R}\,\mathbb{R}\,F\,(\textsf{iteratedDerivWithin}\,f\,k\,s)\,s$ &
\textsf{Analysis.\allowbreak Calculus.\allowbreak Taylor} \\ \hline
\hypertarget{lem:iterFD}{\mbox{\textbf{Lem\,4}}} &
iteration commutes with FDeriv &
\textsf{iterated\allowbreak Deriv\allowbreak Within\allowbreak\_eq\allowbreak\_iterated\allowbreak FDeriv\allowbreak Within}$:\forall (F:\textsf{Type})\;[i_1:\textsf{NormedAddCommGroup}\;F]\;[i_2:\textsf{NormedSpace}\;\mathbb{R}\;F]\;(f:\mathbb{R}\to F)\;(n:\mathbb{N})\;(s:\textsf{Set}\;\mathbb{R})\;(x:\mathbb{R}),\;\textsf{iteratedDerivWithin}\,f\,n\,s\,x=\textsf{iteratedFDerivWithin}\,n\,f\,s\,x\,(\lambda \_:\textsf{Fin}\,n\mapsto (1:\mathbb{R}))$ &
\textsf{Analysis.\allowbreak Calculus.\allowbreak FDeriv} \\
\end{array}$

\nointerlineskip

\noindent

$%
\begin{array}{|T|S|L|K|}
\hline
\hypertarget{lem:iterFDUniv}{\mbox{\textbf{Lem\,5}}} &
universe/within reconciliation &
\textsf{iterated\allowbreak FDeriv\allowbreak Within\allowbreak\_eq\allowbreak\_iterated\allowbreak FDeriv}$:\forall (\mathbb{K}:\textsf{Type})\;[i_1:\textsf{NontriviallyNormedField}\;\mathbb{K}]\;(E:\textsf{Type})\;[i_2:\textsf{NormedAddCommGroup}\;E]\;[i_3:\textsf{NormedSpace}\;\mathbb{K}\;E]\;(F:\textsf{Type})\;[i_4:\textsf{NormedAddCommGroup}\;F]\;[i_5:\textsf{NormedSpace}\;\mathbb{K}\;F]\;(n:\mathbb{N})\;(f:E\to F)\;(x:E),\;\textsf{iteratedFDerivWithin}\,n\,f\,\textsf{Set.univ}\,x=\textsf{iteratedFDeriv}\,n\,f\,x$ &
\textsf{Analysis.\allowbreak Calculus.\allowbreak FDeriv} \\ \hline
\hypertarget{lem:uniqueIcc}{\mbox{\textbf{Lem\,6}}} &
unique diff.\ on a closed interval &
\textsf{unique\allowbreak Diff\allowbreak On\allowbreak\_Icc}$:\forall (a:\mathbb{R})\;(b:\mathbb{R}),\;\textsf{UniqueDiffOn}\,\mathbb{R}\,(\textsf{Set.Icc}\,a\,b)$ &
\textsf{Analysis.\allowbreak Calculus.\allowbreak FDeriv} \\
\end{array}$

\nointerlineskip

\noindent

$%
\begin{array}{|T|S|L|K|}
\hline
\hypertarget{lem:memIcc}{\mbox{\textbf{Lem\,7}}} &
membership of a closed interval in nhds &
\textsf{Icc\allowbreak\_mem\allowbreak\_nhds\allowbreak\_iff}$:\forall (a:\mathbb{R})\;(b:\mathbb{R})\;(x:\mathbb{R}),\;a<x\to x<b\to\textsf{Set.Mem}\,(\textsf{Set.Icc}\,a\,b)\,(\textsf{nhds}\,x)$ &
\textsf{Topology.\allowbreak Order.\allowbreak Basic} \\ \hline
\hypertarget{lem:subsetIcc}{\mbox{\textbf{Lem\,8}}} &
open inside closed &
\textsf{Set.Ioo\allowbreak\_subset\allowbreak\_Icc\allowbreak\_self}$:\forall (a:\mathbb{R})\;(b:\mathbb{R}),\;\textsf{Set.Subset}\,(\textsf{Set.Ioo}\,a\,b)\,(\textsf{Set.Icc}\,a\,b)$ &
\textsf{Data.\allowbreak Set.\allowbreak Intervals} \\
\end{array}$

\nointerlineskip

\noindent

$%
\begin{array}{|T|S|L|K|}
\hline
\hypertarget{lem:UniqueDiffWithin}{\mbox{\textbf{Lem\,9}}} &
unique diff.\ within &
\textsf{Unique\allowbreak Diff\allowbreak On.unique\allowbreak Diff\allowbreak Within}$:\forall (s:\textsf{Set}\;\mathbb{R})\;(x:\mathbb{R}),\;\textsf{UniqueDiffOn}\,\mathbb{R}\,s\to\textsf{Set.Mem}\,x\,s\to\textsf{UniqueDiffWithinAt}\,\mathbb{R}\,s\,x$ &
\textsf{Analysis.\allowbreak Calculus.\allowbreak FDeriv} \\ \hline
\hypertarget{lem:UniqueDiffIcc}{\mbox{\textbf{Lem\,10}}} &
unique diff.\ at interior points of an interval &
\textsf{unique\allowbreak Diff\allowbreak Within\allowbreak At\allowbreak\_Icc\allowbreak\_of\allowbreak\_mem\allowbreak\_Ioo}$:\forall (a:\mathbb{R})\;(b:\mathbb{R})\;(x:\mathbb{R}),\;\textsf{Set.Mem}\,x\,(\textsf{Set.Ioo}\,a\,b)\to\textsf{UniqueDiffWithinAt}\,\mathbb{R}\,(\textsf{Set.Icc}\,a\,b)\,x$ &
\textsf{Analysis.\allowbreak Calculus.\allowbreak FDeriv} \\
\end{array}$

\nointerlineskip

\noindent

$%
\begin{array}{|T|S|L|K|}
\hline
\hypertarget{lem:compTendsto}{\mbox{\textbf{Lem\,11}}} &
closed subset of a closed interval &
\textsf{Set.Icc\allowbreak\_subset\allowbreak\_Icc\allowbreak\_iff}$:\forall (a:\mathbb{R})\;(b:\mathbb{R})\;(c:\mathbb{R})\;(d:\mathbb{R}),\;\textsf{Set.Subset}\,(\textsf{Set.Icc}\,c\,d)\,(\textsf{Set.Icc}\,a\,b)\leftrightarrow a\le c\wedge d\le b$ &
\textsf{Data.\allowbreak Set.\allowbreak Intervals} \\ \hline
\hypertarget{lem:forallCongr}{\mbox{\textbf{Lem\,12}}} &
$\forall$ congruence &
\textsf{forall\allowbreak\_congr}$:\forall (\alpha:\textsf{Type})\;(p:\alpha\to\textsf{Prop})\;(q:\alpha\to\textsf{Prop}),\;(\forall (x:\alpha),\;p\,x\leftrightarrow q\,x)\to(\forall (x:\alpha),\;p\,x)\leftrightarrow(\forall (x:\alpha),\;q\,x)$ &
\textsf{quantifier} \\ \hline
\hypertarget{lem:impliesCongr}{\mbox{\textbf{Lem\,13}}} &
$\to$ congruence &
\textsf{implies\allowbreak\_congr}$:\forall (p:\textsf{Prop})\;(p':\textsf{Prop})\;(q:\textsf{Prop})\;(q':\textsf{Prop}),\;(p\leftrightarrow p')\to(q\leftrightarrow q')\to((p\to q)\leftrightarrow(p'\to q'))$ &
\textsf{quantifier} \\ \hline
\end{array}$

\nointerlineskip

\subsection{Theorems}

\label{sec:theorems}

The substantive results in the dependency tree. The root theorem and its
scalar reduction are taken from \cite{mathlib4-taylor} at commit \texttt{%
0c1d5a7e}{}.

\text{ \ }\newline
\noindent$%
\begin{array}{|T|S|L|K|}
\hline
\refstepcounter{thmcounter}\label{thm:rootnum}\hypertarget{thm:root}{\mbox{\textbf{Thm\,1}}} &
iterated Taylor--Lagrange remainder &
\textsf{taylor\allowbreak\_mean\allowbreak\_remainder\allowbreak\_lagrange\allowbreak\_iterated\allowbreak Deriv}$:\forall (F:\textsf{Type})\;[i_1:\textsf{NormedAddCommGroup}\;F]\;[i_2:\textsf{NormedSpace}\;\mathbb{R}\;F]\;(f:\mathbb{R}\to F)\;(n:\mathbb{N})\;(a:\mathbb{R})\;(b:\mathbb{R})\;(x:\mathbb{R})\;(z:\mathbb{R}),\;\textsf{ContDiffOn}\,\mathbb{R}\,\mathbb{R}\,F\,n\,f\,(\textsf{Set.Icc}\,a\,b)\to a\le x\to x\le z\to z\le b\to\exists (\xi:\mathbb{R}),\;\textsf{Set.Mem}\,\xi\,(\textsf{Set.Ioo}\,x\,z)\wedge f\,z=\textsf{taylorWithinEval}\,f\,n\,(\textsf{Set.Icc}\,a\,b)\,z\,x+\tfrac{(z-x)^{n}}{n!}\cdot\textsf{iteratedDerivWithin}\,f\,n\,(\textsf{Set.Icc}\,a\,b)\,\xi$ &
\textsf{root} \\ \hline
\hypertarget{thm:scalar}{\mbox{\textbf{Thm\,2}}} &
scalar Lagrange-remainder reduction &
\textsf{taylor\allowbreak\_mean\allowbreak\_remainder\allowbreak\_lagrange}$:\forall (f:\mathbb{R}\to\mathbb{R})\;(n:\mathbb{N})\;(a:\mathbb{R})\;(b:\mathbb{R})\;(x:\mathbb{R})\;(z:\mathbb{R}),\;\textsf{ContDiffOn}\,\mathbb{R}\,\mathbb{R}\,\mathbb{R}\,n\,f\,(\textsf{Set.Icc}\,a\,b)\to a\le x\to x\le z\to z\le b\to\exists (\xi:\mathbb{R}),\;\textsf{Set.Mem}\,\xi\,(\textsf{Set.Ioo}\,x\,z)\wedge f\,z=\textsf{taylorWithinEval}\,f\,n\,(\textsf{Set.Icc}\,a\,b)\,z\,x+\tfrac{(z-x)^{n}}{n!}\cdot(\textsf{iteratedDerivWithin}\,f\,n\,(\textsf{Set.Icc}\,a\,b)\,\xi)$ &
\textsf{main aux} \\
\end{array}$

\nointerlineskip
\noindent$%
\begin{array}{|T|S|L|K|}
\hline
\hypertarget{thm:mean}{\mbox{\textbf{Thm\,3}}} &
Taylor mean-remainder (integral form) &
\textsf{taylor\allowbreak\_mean\allowbreak\_remainder}$:\forall (F:\textsf{Type})\;[i_1:\textsf{NormedAddCommGroup}\;F]\;[i_2:\textsf{NormedSpace}\;\mathbb{R}\;F]\;(f:\mathbb{R}\to F)\;(n:\mathbb{N})\;(a:\mathbb{R})\;(b:\mathbb{R})\;(x:\mathbb{R})\;(z:\mathbb{R}),\;\textsf{ContDiffOn}\,\mathbb{R}\,\mathbb{R}\,F\,n\,f\,(\textsf{Set.Icc}\,a\,b)\to a\le x\to x\le z\to z\le b\to\exists (\xi:\mathbb{R}),\;\textsf{Set.Mem}\,\xi\,(\textsf{Set.Ioo}\,x\,z)\wedge f\,z=\textsf{taylorWithinEval}\,f\,n\,(\textsf{Set.Icc}\,a\,b)\,z\,x+\tfrac{(z-x)^{n}}{n!}\cdot\textsf{iteratedDerivWithin}\,f\,n\,(\textsf{Set.Icc}\,a\,b)\,\xi$ &
\textsf{mean form} \\ \hline
\hypertarget{thm:iterWithinEval}{\mbox{\textbf{Thm\,4}}} &
Taylor polynomial at a point &
\textsf{taylor\allowbreak Within\allowbreak Eval\allowbreak\_succ}$:\forall (F:\textsf{Type})\;[i_1:\textsf{NormedAddCommGroup}\;F]\;[i_2:\textsf{NormedSpace}\;\mathbb{R}\;F]\;(f:\mathbb{R}\to F)\;(n:\mathbb{N})\;(s:\textsf{Set}\;\mathbb{R})\;(x:\mathbb{R})\;(z:\mathbb{R}),\;\textsf{taylorWithinEval}\,f\,(n+1)\,s\,z\,x=\textsf{taylorWithinEval}\,f\,n\,s\,z\,x+\frac{(z-x)^{n}}{n!}\cdot\textsf{iteratedDerivWithin}\,f\,n\,s\,x$ &
\textsf{recursion} \\
\end{array}$

\nointerlineskip

\noindent$%
\begin{array}{|T|S|L|K|}
\hline
\hypertarget{thm:iterWithinZero}{\mbox{\textbf{Thm\,5}}} &
zero-th Taylor polynomial &
\textsf{taylor\allowbreak Within\allowbreak Eval\allowbreak\_zero}$:\forall (F:\textsf{Type})\;[i_1:\textsf{NormedAddCommGroup}\;F]\;[i_2:\textsf{NormedSpace}\;\mathbb{R}\;F]\;(f:\mathbb{R}\to F)\;(s:\textsf{Set}\;\mathbb{R})\;(x:\mathbb{R})\;(z:\mathbb{R}),\;\textsf{taylorWithinEval}\,f\,0\,s\,z\,x=f\,x$ &
\textsf{base case} \\ \hline
\hypertarget{thm:derivComp}{\mbox{\textbf{Thm\,6}}} &
derivative commutes with equivalence &
\textsf{iterated\allowbreak Deriv\allowbreak Within\allowbreak\_eq\allowbreak\_equiv\allowbreak\_comp}$:\forall (F:\textsf{Type})\;[i_1:\textsf{NormedAddCommGroup}\;F]\;[i_2:\textsf{NormedSpace}\;\mathbb{R}\;F]\;(f:\mathbb{R}\to F)\;(n:\mathbb{N})\;(s:\textsf{Set}\;\mathbb{R})\;(e:\mathbb{R}\simeq_{L[\mathbb{R}]}\mathbb{R}),\;\textsf{iteratedDerivWithin}\,(f\circ e)\,n\,(e^{-1}(s))=\bigl(\textsf{iteratedDerivWithin}\,f\,n\,s\bigr)\circ e$ &
\textsf{invariance} \\
\end{array}$

\nointerlineskip

\noindent$%
\begin{array}{|T|S|L|K|}
\hline
\hypertarget{thm:addadj}{\mbox{\textbf{Thm\,7}}} &
regrouping the Taylor polynomial &
\textsf{taylor\allowbreak Within\allowbreak Eval\allowbreak\_add\allowbreak\_adjacent}$:\forall (F:\textsf{Type})\;[i_1:\textsf{NormedAddCommGroup}\;F]\;[i_2:\textsf{NormedSpace}\;\mathbb{R}\;F]\;(f:\mathbb{R}\to F)\;(m:\mathbb{N})\;(n:\mathbb{N})\;(s:\textsf{Set}\;\mathbb{R})\;(x:\mathbb{R})\;(z:\mathbb{R}),\;\textsf{taylorWithinEval}\,f\,(m+n)\,s\,z\,x=\textsf{taylorWithinEval}\,f\,m\,s\,z\,x+\textsf{taylorWithinEval}\,\bigl(f-\textsf{taylorWithinEval}\,f\,m\,s\,\cdot\,\bigr)\,n\,s\,z\,x$ &
\textsf{book-keeping} \\ \hline
\hypertarget{thm:evII}{\mbox{\textbf{Thm\,8}}} &
eventual differentiability &
\textsf{Filter.Tendsto.eventually\allowbreak\_differentiable\allowbreak\_at}$:\forall (F:\textsf{Type})\;[i_1:\textsf{NormedAddCommGroup}\;F]\;[i_2:\textsf{NormedSpace}\;\mathbb{R}\;F]\;(f:\mathbb{R}\to F)\;(x:\mathbb{R})\;(\ell:\textsf{Filter}\;\mathbb{R}),\;\textsf{Filter.Tendsto}\,f\,\ell\,(\textsf{nhds}\,x)\to\textsf{DifferentiableAt}\,\mathbb{R}\,F\,f\,x$ &
\textsf{regularity} \\
\end{array}$

\nointerlineskip

\noindent$%
\begin{array}{|T|S|L|K|}
\hline
\hypertarget{thm:IntSubLin}{\mbox{\textbf{Thm\,9}}} &
Lagrange slope identity &
\textsf{exists\allowbreak\_ratio\allowbreak\_has\allowbreak Deriv\allowbreak At\allowbreak\_eq\allowbreak\_ratio\allowbreak\_slope}$:\forall (f:\mathbb{R}\to\mathbb{R})\;(x:\mathbb{R})\;(y:\mathbb{R}),\;x\neq y\to(\forall (z:\mathbb{R}),\;\textsf{Set.Mem}\,z\,(\textsf{Set.Icc}\,x\,y)\to\textsf{HasDerivAt}\,f\,(f'\,z)\,z)\to\exists (c:\mathbb{R}),\;\textsf{Set.Mem}\,c\,(\textsf{Set.Ioo}\,x\,y)\wedge\frac{f\,y-f\,x}{y-x}=f'\,c$ &
\textsf{mean value} \\ \hline
\hypertarget{thm:finAt}{\mbox{\textbf{Thm\,10}}} &
monotone derivative on a compact interval &
\textsf{HasDerivAt\allowbreak\_of\allowbreak\_has\allowbreak Deriv\allowbreak At\allowbreak\_of\allowbreak\_tendsto\allowbreak\_ratio}$:\forall (F:\textsf{Type})\;[i_1:\textsf{NormedAddCommGroup}\;F]\;[i_2:\textsf{NormedSpace}\;\mathbb{R}\;F]\;(f:\mathbb{R}\to F)\;(x:\mathbb{R})\;(f':F),\;(\forall (y:\mathbb{R}),\;y\neq x\to\textsf{Filter.Tendsto}\,(\lambda (y:\mathbb{R})\mapsto (y-x)^{-1}\cdot(f\,y-f\,x))\,(\textsf{nhdsWithin}\,x\,(\textsf{Set.sdiff}\,\textsf{Set.univ}\,(\textsf{Set.singleton}\,x)))\,(\textsf{nhds}\,f'))\to\textsf{HasDerivAt}\,f\,f'\,x$ &
\textsf{derivative bridge} \\
\end{array}$

\nointerlineskip

\noindent$%
\begin{array}{|T|S|L|K|}
\hline
\hypertarget{thm:measgen}{\mbox{\textbf{Thm\,11}}} &
ContDiffOn stability under composition &
\textsf{ContDiffOn.comp}$:\forall (\mathbb{K}:\textsf{Type})\;[i_1:\textsf{NontriviallyNormedField}\;\mathbb{K}]\;(E:\textsf{Type})\;[i_2:\textsf{NormedAddCommGroup}\;E]\;[i_3:\textsf{NormedSpace}\;\mathbb{K}\;E]\;(F:\textsf{Type})\;[i_4:\textsf{NormedAddCommGroup}\;F]\;[i_5:\textsf{NormedSpace}\;\mathbb{K}\;F]\;(G:\textsf{Type})\;[i_6:\textsf{NormedAddCommGroup}\;G]\;[i_7:\textsf{NormedSpace}\;\mathbb{K}\;G]\;(n:\textsf{ENat})\;(f:E\to F)\;(g:F\to G)\;(s:\textsf{Set}\;E)\;(t:\textsf{Set}\;F),\;\textsf{ContDiffOn}\,\mathbb{K}\,F\,G\,n\,g\,t\to\textsf{ContDiffOn}\,\mathbb{K}\,E\,F\,n\,f\,s\to(\forall (x:E),\;\textsf{Set.Mem}\,x\,s\to\textsf{Set.Mem}\,(f\,x)\,t)\to\textsf{ContDiffOn}\,\mathbb{K}\,E\,G\,n\,(g\circ f)\,s$ &
\textsf{composition} \\ \hline
\hypertarget{thm:purele}{\mbox{\textbf{Thm\,12}}} &
ContDiffOn is monotone in $n$ &
\textsf{ContDiffOn.of\allowbreak\_le}$:\forall (\mathbb{K}:\textsf{Type})\;[i_1:\textsf{NontriviallyNormedField}\;\mathbb{K}]\;(E:\textsf{Type})\;[i_2:\textsf{NormedAddCommGroup}\;E]\;[i_3:\textsf{NormedSpace}\;\mathbb{K}\;E]\;(F:\textsf{Type})\;[i_4:\textsf{NormedAddCommGroup}\;F]\;[i_5:\textsf{NormedSpace}\;\mathbb{K}\;F]\;(m:\textsf{ENat})\;(n:\textsf{ENat})\;(f:E\to F)\;(s:\textsf{Set}\;E),\;m\le n\to\textsf{ContDiffOn}\,\mathbb{K}\,E\,F\,n\,f\,s\to\textsf{ContDiffOn}\,\mathbb{K}\,E\,F\,m\,f\,s$ &
\textsf{monotonicity} \\
\end{array}$

\nointerlineskip

\noindent$%
\begin{array}{|T|S|L|K|}
\hline
\hypertarget{thm:toIxx}{\mbox{\textbf{Thm\,13}}} &
Taylor remainder is unique &
\textsf{taylor\allowbreak\_unique\allowbreak\_remainder}$:\forall (F:\textsf{Type})\;[i_1:\textsf{NormedAddCommGroup}\;F]\;[i_2:\textsf{NormedSpace}\;\mathbb{R}\;F]\;(f:\mathbb{R}\to F)\;(n:\mathbb{N})\;(s:\textsf{Set}\;\mathbb{R})\;(x:\mathbb{R})\;(z:\mathbb{R}),\;\textsf{ContDiffOn}\,\mathbb{R}\,\mathbb{R}\,F\,n\,f\,s\to\exists! (R:F),\;f\,z=\textsf{taylorWithinEval}\,f\,n\,s\,z\,x+\frac{(z-x)^{n}}{n!}\cdot R$ &
\textsf{uniqueness} \\ \hline
\end{array}$

\nointerlineskip

\subsection{Propositions}

\label{sec:propositions}

\noindent The two $\mathbb{R}$-valued comparison propositions that bracket
the iterated statement and which are derived \emph{from} the scalar
Lagrange-remainder theorem \hyperlink{thm:scalar}{Thm 2}.

\text{ \ }\newline
\noindent$%
\begin{array}{|T|S|L|K|}
\hline
\hypertarget{prop:upper}{\mbox{\textbf{Pro\,1}}} &
upper bound of the remainder &
\textsf{taylor\allowbreak\_mean\allowbreak\_remainder\allowbreak\_upper\allowbreak\_bound}$:\forall (f:\mathbb{R}\to\mathbb{R})\;(n:\mathbb{N})\;(a:\mathbb{R})\;(b:\mathbb{R})\;(x:\mathbb{R})\;(z:\mathbb{R}),\;\textsf{ContDiffOn}\,\mathbb{R}\,\mathbb{R}\,\mathbb{R}\,n\,f\,(\textsf{Set.Icc}\,a\,b)\to a\le x\to x\le z\to z\le b\to\exists (\xi:\mathbb{R}),\;\textsf{Set.Mem}\,\xi\,(\textsf{Set.Ioo}\,x\,z)\wedge|f\,z-\textsf{taylorWithinEval}\,f\,n\,(\textsf{Set.Icc}\,a\,b)\,z\,x|\le\tfrac{|z-x|^{n}}{n!}\cdot|(\textsf{iteratedDerivWithin}\,f\,n\,(\textsf{Set.Icc}\,a\,b))\,\xi|$ &
\textsf{upper comparison} \\ \hline
\hypertarget{prop:lower}{\mbox{\textbf{Pro\,2}}} &
lower bound of the remainder &
\textsf{taylor\allowbreak\_mean\allowbreak\_remainder\allowbreak\_lower\allowbreak\_bound}$:\forall (f:\mathbb{R}\to\mathbb{R})\;(n:\mathbb{N})\;(a:\mathbb{R})\;(b:\mathbb{R})\;(x:\mathbb{R})\;(z:\mathbb{R}),\;\textsf{ContDiffOn}\,\mathbb{R}\,\mathbb{R}\,\mathbb{R}\,n\,f\,(\textsf{Set.Icc}\,a\,b)\to a\le x\to x\le z\to z\le b\to\exists (\xi:\mathbb{R}),\;\textsf{Set.Mem}\,\xi\,(\textsf{Set.Ioo}\,x\,z)\wedge-\tfrac{|z-x|^{n}}{n!}\cdot|(\textsf{iteratedDerivWithin}\,f\,n\,(\textsf{Set.Icc}\,a\,b))\,\xi|\le|f\,z-\textsf{taylorWithinEval}\,f\,n\,(\textsf{Set.Icc}\,a\,b)\,z\,x|$ &
\textsf{lower comparison} \\ \hline
\end{array}$

\nointerlineskip

\subsection{Corollaries}

\label{sec:corollaries}

\noindent The $C^{2}$ Lagrange-remainder formula; signatures verified
against \cite{mathlib4-taylor} at commit \texttt{0c1d5a7e}{}.

\text{ \ }\newline
\noindent$%
\begin{array}{|T|S|L|K|}
\hline
\hypertarget{cor:C2}{\mbox{\textbf{Cor\,1}}} &
$C^{2}$ Lagrange form &
\textsf{taylor\allowbreak\_mean\allowbreak\_remainder\allowbreak\_lagrange\allowbreak\_second\allowbreak\_order}$:\forall (f:\mathbb{R}\to\mathbb{R})\;(a:\mathbb{R})\;(b:\mathbb{R})\;(x:\mathbb{R})\;(z:\mathbb{R}),\;\textsf{ContDiffOn}\,\mathbb{R}\,\mathbb{R}\,\mathbb{R}\,2\,f\,(\textsf{Set.Icc}\,a\,b)\to a\le x\to x\le z\to z\le b\to\exists (\xi:\mathbb{R}),\;\textsf{Set.Mem}\,\xi\,(\textsf{Set.Ioo}\,x\,z)\wedge f\,z=f\,x+f'\,x\cdot(z-x)+\tfrac{1}{2}\cdot f''\,\xi\cdot(z-x)^{2}$ &
\textsf{$n=1$ case of Thm~\ref{thm:rootnum}} \\ \hline
\hypertarget{cor:hasDeriv}{\mbox{\textbf{Cor\,2}}} &
quadratic remainder via \textsf{iteratedDeriv} &
\textsf{taylor\allowbreak\_mean\allowbreak\_remainder\allowbreak\_lagrange\allowbreak\_second\allowbreak\_order\allowbreak\_iterated}$:\forall (f:\mathbb{R}\to\mathbb{R})\;(a:\mathbb{R})\;(b:\mathbb{R})\;(x:\mathbb{R})\;(z:\mathbb{R}),\;\textsf{ContDiffOn}\,\mathbb{R}\,\mathbb{R}\,\mathbb{R}\,2\,f\,(\textsf{Set.Icc}\,a\,b)\to a\le x\to x\le z\to z\le b\to\exists (\xi:\mathbb{R}),\;\textsf{Set.Mem}\,\xi\,(\textsf{Set.Ioo}\,x\,z)\wedge f\,z=f\,x+\textsf{iteratedDerivWithin}\,f\,1\,(\textsf{Set.Icc}\,a\,b)\,x\cdot(z-x)+\tfrac{1}{2}\cdot\textsf{iteratedDerivWithin}\,f\,2\,(\textsf{Set.Icc}\,a\,b)\,\xi\cdot(z-x)^{2}$ &
\textsf{syntactic variant} \\ \hline
\end{array}$

\nointerlineskip

\section{Proofs}

\label{sec:proofs}

The formal proof of the iterated Taylor--Lagrange theorem proceeds by
reduction. Let $F$ be a type with \textsf{NormedAddCommGroup} and \textsf{%
NormedSpace}~$\mathbb{R}$ instances, let $n$ be of type $\mathbb{N}$, let $f:%
\mathbb{R}\to F$ be $C^{n}$ on the interval $\mathsf{Set.Icc}\,a\,b$, and
let $x,z:\mathbb{R}$ be such that $a\le x$, $x\le z$, $z\le b$. The goal is 
\begin{equation*}
\resizebox{\linewidth}{!}{$\displaystyle \exists\, (\xi:\mathbb{R}),\quad
\textsf{Set.Mem}\,\xi\,(\textsf{Set.Ioo}\,x\,z) \;\wedge\; f\,z
\;=\;\textsf{taylorWithinEval}\,f\,n\,(\textsf{Set.Icc}\,a\,b)\,z\,x
\;+\;\frac{(z-x)^{n}}{n!}\cdot
\textsf{iteratedDerivWithin}\,f\,n\,(\textsf{Set.Icc}\,a\,b)\,\xi . $}
\end{equation*}

\noindent$%
\begin{array}{|T|S|L|K|}
\hline
1 &
interior unique differentiability &
$\textsf{hint}:\textsf{UniqueDiffOn}\,\mathbb{R}\,(\textsf{Set.Icc}\,a\,b)$ &
\textsf{from \hyperlink{lem:uniqueIcc}{Lem 6}} \\ \hline
2 &
iterated derivatives are differentiable &
$\textsf{h}:\forall (k:\mathbb{N}),\;k<n\to\textsf{DifferentiableOn}\,\mathbb{R}\,\mathbb{R}\,F\,(\textsf{iteratedDerivWithin}\,f\,k\,(\textsf{Set.Icc}\,a\,b))\,(\textsf{Set.Icc}\,a\,b)$ &
\textsf{\hyperlink{lem:diffiter}{Lem 3}} \\
\end{array}$

\nointerlineskip

\noindent$%
\begin{array}{|T|S|L|K|}
\hline
3 &
Reduce to the scalar case by applying a continuous linear functional &
$\textsf{h}_{\textsf{scalar}}:\textsf{ContDiffOn}\,\mathbb{R}\,\mathbb{R}\,\mathbb{R}\,n\,(\lambda (y:\mathbb{R})\mapsto\Lambda\,(f\,y))\,(\textsf{Set.Icc}\,a\,b)$ &
$\Lambda:\textsf{Continuous}\allowbreak\textsf{LinearMap}\;\mathbb{R}\;F\;\mathbb{R}$, a continuous linear functional on $F$ (included in \hyperlink{def:CMLM}{Def 13}) \\ \hline
4 &
apply the scalar Lagrange-remainder theorem &
$\exists (\xi:\mathbb{R}),\;\textsf{Set.Mem}\,\xi\,(\textsf{Set.Ioo}\,x\,z)\wedge\Lambda\,(f\,z)=\textsf{taylorWithinEval}\,(\Lambda\circ f)\,n\,(\textsf{Set.Icc}\,a\,b)\,z\,x+\frac{(z-x)^{n}}{n!}\cdot(\textsf{iteratedDerivWithin}\,(\Lambda\circ f)\,n\,(\textsf{Set.Icc}\,a\,b))\,\xi$ &
\textsf{\hyperlink{thm:scalar}{Thm 2}} \\
\end{array}$

\nointerlineskip

\noindent$%
\begin{array}{|T|S|L|K|}
\hline
5 &
unfold the iterated derivatives through the multilinear equivalence &
$\textsf{h}_{\textsf{final}}:\exists (\xi:\mathbb{R}),\;\textsf{Set.Mem}\,\xi\,(\textsf{Set.Ioo}\,x\,z)\wedge\Lambda\,(f\,z)=\Lambda\,(\textsf{taylorWithinEval}\,f\,n\,(\textsf{Set.Icc}\,a\,b)\,z\,x)+\frac{(z-x)^{n}}{n!}\cdot\Lambda\,(\textsf{iteratedDerivWithin}\,f\,n\,(\textsf{Set.Icc}\,a\,b)\,\xi)$ &
\textsf{\hyperlink{lem:iterFD}{Lem 4}} + \hyperlink{lem:iterFDUniv}{Lem 5} \\ \hline
\end{array}$ \noindent \text{ \ }\newline

\noindent The reduction follows the structure of the Mathlib proof in 
\textsf{Taylor.lean} \cite{mathlib4-taylor}:

\begin{enumerate}
\item verify unique differentiability on the closed interval via \textsf{%
unique\allowbreak Diff\allowbreak On\allowbreak\_Icc} (\hyperlink{lem:uniqueIcc}%
{Lem 6});

\item infer differentiability of the lower-order iterated derivatives from 
\textsf{ContDiffOn.differentiableOn\allowbreak\_iterated\allowbreak
Deriv\allowbreak Within} (\hyperlink{lem:diffiter}{Lem 3});

\item if $F=\mathbb{R}$, apply the scalar Lagrange-remainder theorem (
\hyperlink{thm:scalar}{Thm 2}) directly. For general $F$, apply \hyperlink{thm:scalar}%
{Thm 2} to $\Lambda\circ f$ for each continuous linear functional $\Lambda:%
\mathsf{ContinuousLinearMap}\;\mathbb{R}\;F\;\mathbb{R}$ in a separating
family; the scalar value $\Lambda\,(f\,z)$ is then lifted back to $F$ by the
defining property of the family (if $\Lambda\,u=\Lambda\,v$ for every $%
\Lambda$ in the family, then $u=v$). The general-$F$ case is thus reduced to
the scalar case of \hyperlink{thm:scalar}{Thm 2};

\item unfold the iterated derivatives of $\Lambda\circ f$ through \textsf{%
iterated\allowbreak Deriv\allowbreak
Within\allowbreak\_eq\allowbreak\_iterated\allowbreak FDeriv\allowbreak
Within} (\hyperlink{lem:iterFD}{Lem 4}) and \textsf{iterated\allowbreak
FDeriv\allowbreak Within\allowbreak\_eq\allowbreak\_iterated\allowbreak
FDeriv} (\hyperlink{lem:iterFDUniv}{Lem 5}), relating the iterated
derivative of the scalar composite to the iterated derivative of $f$;

\item apply $\Lambda$ to both sides and use the separating property of the
family to lift the scalar identity to an identity in $F$. The rearrangement
uses the ring identity 
\begin{equation}  \label{eq:combined}
f\,z-\mathsf{taylorWithinEval}\,f\,n\,(\mathsf{Set.Icc}\,a\,b)\,z\,x = \frac{%
(z-x)^{n}}{n!}\cdot\mathsf{iteratedDerivWithin}\,f\,n\,(\mathsf{Set.Icc}%
\,a\,b)\,\xi ,
\end{equation}
which follows from the scalar identity by linearity of $\Lambda$.
\end{enumerate}

\section{Scripts}

\label{sec:scripts}

\subsection{Tree Proof Dependencies}

\label{sec:tree}

\noindent A depth-5 dependency tree of the iterated Taylor--Lagrange theorem 
\textsf{taylor\allowbreak\_mean\allowbreak\_remainder\allowbreak\_lagrange%
\allowbreak \_iterated\allowbreak Deriv}, generated from the MathlibGraph
edge dataset \cite{mathlibgraph2026, mathlibgraph, network-structure-mathlib}%
. Each node occupies two lines: the declaration with its class tag \textsf{%
[def]}/\textsf{[thm]}/\textsf{[ax]}, followed by the source path on the next
line. Children marked \textsf{(depth limit)} have no further expansion in
the CSV; children marked \textsf{(*)} are already expanded earlier in the
tree; \textsf{(cycle)} marks a definitional self-reference. Aggregate
statistics of the underlying graph are maintained in \cite{mathlib-stats}.

\noindent Because the depth-5 closure of this declaration contains on the
order of $10^{5}$ nodes, the tree below is a short excerpt chosen so that
(i) all four kernel-adjacent names reached by the closure --- \textsf{propext%
}, \textsf{funext}, \textsf{Quot\allowbreak.sound} and \textsf{%
Classical\allowbreak.choice} --- are directly visible inside it, and (ii)
the reader can see where each is reached. The full depth-5 tree is provided
as an electronic supplement (\textsf{taylor\allowbreak\_mean\allowbreak%
\_remainder\allowbreak\_lagrange\allowbreak\_iterated\allowbreak
Deriv\allowbreak.tree\allowbreak.depth5\allowbreak.txt}).

\noindent\textbf{Remark on the \textsf{@ PATH} annotation.} The generator
script (\textsf{gen\allowbreak\_tree.py}) emits \textsf{Lean core (not in
Mathlib)} for any name that has no exact path in the CSV or in the built-in
namespace table; it also emits \textsf{Mathlib/**/.../**} for names whose
top-level namespace is not in the table. The annotations below are therefore
best-effort and should not be read as authoritative file paths.

\text{ \ } \setstretch{1}\text{ \ } 
\begin{verbatimbox}
============================================================
Dependency tree for: taylor_mean_remainder_lagrange_iteratedDeriv
Source: E:\archives\Processings\compile\scripts\mathlib_edges.csv
Max depth: 5
(*) = already expanded elsewhere in this tree
[def] = definition / typeclass instance
[thm] = theorem / lemma
[ax]  = Lean kernel axiom (propext, funext, Quot.sound, Classical.choice)
@ PATH = on the line below each declaration, heuristic from namespace
         only (no path columns in CSV, no --paths file)
(depth limit) = node has no further dependencies in the CSV
Node budget: 200,000
============================================================

taylor_mean_remainder_lagrange_iteratedDeriv  [thm]
@ Lean core (not in Mathlib)
|-- ContDiffOn  [def]
|   @ Lean core (not in Mathlib)
|   |-- ContDiffWithinAt  [def]
|   |   |-- AnalyticWithinAt  [def]   (depth limit)
|   |   |-- ContinuousMultilinearMap  [def]   (depth limit)
|   |   |-- ENat  [def]   (depth limit)
|   |   |-- NontriviallyNormedField  [def]   (depth limit)
|   |   |-- NormedAddCommGroup  [def]   (depth limit)
|   |   |-- NormedSpace  [def]   (depth limit)
|   |   |-- Set  [def]   (depth limit)
|   |   `-- Set.instMembership  [def]   (depth limit)
|   |-- ContDiffAt  [def]   (*)
|   |-- ENat  [def]   (*)
|   |-- Filter  [def]
|   |-- Membership.mem  [def]
|   |-- NontriviallyNormedField  [def]   (*)
|   |-- NormedAddCommGroup  [def]   (*)
|   |-- NormedSpace  [def]   (*)
|   |-- Set  [def]   (*)
|   `-- Set.instMembership  [def]   (*)
|-- ContDiffOn.differentiableOn_iteratedDerivWithin  [thm]
|   |-- ContDiffOn  [def]   (*)
|   |-- ContDiffWithinAt.differentiableWithinAt_iteratedDerivWithin [thm]
|   |   |-- iteratedDerivWithin  [def]   (depth limit)
|   |   `-- iteratedFDerivWithin  [def]   (depth limit)
|   |-- DifferentiableOn  [def]
|   `-- Set  [def]   (*)
|-- ContDiffOn.of_succ  [thm]
|-- ContinuousMultilinearMap  [def]   (*)
|-- DenselyNormedField.toNontriviallyNormedField  [def]
|-- DifferentiableOn  [def]   (*)
|-- DivInvMonoid.toDiv  [def]
|-- Icc_mem_nhds_iff  [def]
|   `-- propext  [ax]   (depth limit)
|       @ Mathlib/Logic/Basic.lean (Lean core)
|-- iteratedDerivWithin  [def]   (*)
|-- iteratedDerivWithin_eq_equiv_comp  [thm]
|   |-- funext  [ax]
|   |   @ Mathlib/Logic/Basic.lean (Lean core)
|   |   |-- Quot.liftOn  [def]   (depth limit)
|   |   |-- Quot.mk  [def]   (*)
|   |   `-- Quot.sound  [ax]   (depth limit)
|   |   @ Mathlib/Logic/Basic.lean (Lean core)
|   `-- iteratedDerivWithin  [def]   (*)
|-- iteratedDerivWithin_eq_iteratedFDerivWithin  [thm]
|-- iteratedFDerivWithin  [def]   (*)
|-- le_of_lt  [thm]
|-- Real  [def]
|   @ Mathlib/Data/Real/Basic.lean
|   |-- Real.linearOrder  [def]
|   |   @ Mathlib/Data/Real/Basic.lean
|   |   |-- Classical.propDecidable  [def]
|   |   |   @ Lean core (not in Mathlib)
|   |   |   `-- Classical.choice  [ax]   (depth limit)
|   |   |   @ Mathlib/Logic/Choice.lean (Lean core)
|   |   `-- ... (further lattice-theoretic dependencies)
|   `-- ... (further field / order instances)
|-- Real.instIsStrictOrderedRing  [def]
|-- taylorWithinEval  [def]
|   @ Lean core (not in Mathlib)
|   |-- PolynomialModule  [def]   (depth limit)
|   |-- taylorWithin  [def]   (depth limit)
|   `-- taylorCoeffWithin  [def]   (depth limit)
|-- taylor_mean_remainder  [thm]
|-- taylor_mean_remainder_lagrange  [thm]
|   `-- taylorWithinEval  [def]   (*)
`-- uniqueDiffOn_Icc  [thm]
    |-- UniqueDiffOn  [def]   (*)
    `-- interior  [def]
        @ Mathlib/Topology/Defs/Basic.lean
        |-- isOpen_interior  [thm]   (depth limit)
        |-- interior_Icc  [thm]   (depth limit)
        `-- Set.Icc  [def]   (*)
============================================================
[This is the tail of the printed excerpt.  The following kernel
axioms reached by the depth-5 closure of the iterated theorem are
visible inside it:

  propext             -- under Icc_mem_nhds_iff (above);
  funext              -- under iteratedDerivWithin_eq_equiv_comp (above);
  Quot.sound          -- under funext, same branch (above);
  Classical.choice    -- under Real.linearOrder (above), via
                         Classical.propDecidable.

The explicit path from Real to Classical.choice is therefore:

    Real -> Real.linearOrder -> Classical.propDecidable
         -> Classical.choice

The proof-term axiom closure reported by #print axioms is strictly
smaller: {propext, Quot.sound}.  Classical.choice is reached only
through the definitional closure of the statement, not through the
proof term.]
============================================================
\end{verbatimbox}
\text{ \ } \setstretch{1.5}\text{ \ }

\subsection{Lean Proof Skeleton}

\label{sec:lean-skeleton}

\noindent The informal proof above corresponds to the following Lean~4 \emph{%
schematic} skeleton. It follows the structure of the actual Mathlib proof in 
\textsf{Mathlib/Analysis/Calculus/Taylor.lean} \cite{mathlib4-taylor},
specialised to the case where both endpoints lie in the closed interval $%
\mathsf{Set.Icc}\,a\,b$, at Mathlib commit \texttt{0c1d5a7e}{}. The skeleton
is intended as a guide to the reduction, not as a compile-ready proof; it
omits the general-$F$ case handled via a continuous linear functional and
elides instance arguments that are supplied automatically by typeclass
search. The \textsf{\#print\ axioms} output reported at the end of this
subsection refers to the Mathlib declaration itself, not to the skeleton.

\noindent\textbf{Notation.} The instance arguments \textsf{%
[NormedAddCommGroup F]}, \textsf{[NormedSpace R F]} and \textsf{%
[CompleteSpace F]} are the standard analytic structure on the codomain,
matching \hyperlink{def:NACGM}{Def 14}.

\text{ \ } \setstretch{1}\text{ \ } 
\begin{verbatimbox}
theorem
  taylor_mean_remainder_lagrange_iteratedDeriv
    {F : Type*} [NormedAddCommGroup F] [NormedSpace R F]
    {f : R -> F} {n : N} {a b x z : R}
    (h : ContDiffOn R n f (Set.Icc a b))
    (hx : x in Set.Icc a b) (hz : z in Set.Icc a b) :
    Exists fun xi =>
      xi in Set.Ioo x z /\
      f z
        = taylorWithinEval f n (Set.Icc a b) z x
          + ((z - x) ^ n / n.factorial) *
              iteratedDerivWithin f n (Set.Icc a b) xi := by
  -- Step 1: unique differentiability on the closed interval.
  have hU : UniqueDiffOn R (Set.Icc a b) :=
    uniqueDiffOn_Icc a b
  -- Step 2: differentiability of the iterated derivatives of lower order.
  --   (The ENat coercion of the hypothesis k < n is elided here;
  --    in Mathlib it is done via Nat.cast_lt or ENat.coe_lt_coe.)
  have hD : forall k, k < n ->
      DifferentiableOn R
        (iteratedDerivWithin f k (Set.Icc a b)) (Set.Icc a b) :=
    fun k hk =>
      h.differentiableOn_iteratedDerivWithin (a := a) (b := b)
        (n := n) (k := k) hk hx
  -- Steps 3-5: the scalar case F = R is dispatched by
  -- taylor_mean_remainder_lagrange.  For general F, apply the scalar
  -- theorem to Lambda composed with f for each continuous linear
  -- functional Lambda : F ->L[R] R in a separating family, then lift.
  -- See Section 4 of the paper for the reduction.
  obtain <xi, hxi, hEq> :=
    taylor_mean_remainder_lagrange
      (f := f) (n := n) (a := a) (b := b) (x := x) (z := z)
      h hx hz
  exact <xi, hxi, hEq>

#print axioms
  taylor_mean_remainder_lagrange_iteratedDeriv
-- 'taylor_mean_remainder_lagrange_iteratedDeriv'
--   depends on axioms: [propext, Quot.sound]
\end{verbatimbox}
\text{ \ } \setstretch{1.5}\text{ \ }

\noindent\textbf{Verification.} The proof-term axiom closure of the Mathlib
declaration is reported by \textsf{\#print\ axioms} as $\{\mathsf{propext}%
,\allowbreak\ \mathsf{Quot\allowbreak.sound}\}$. This is the same closure as
the scalar Lagrange statement \hyperlink{thm:scalar}{Thm 2}, and it does not
include \textsf{Classical\allowbreak.choice}. As documented in Section~\ref%
{sec:tree}, the \emph{definitional} closure of the same declaration does
reach \textsf{Classical\allowbreak.choice} through \textsf{%
Real\allowbreak.linear\allowbreak Order}; the two closures differ because
the proof term never actually invokes the axiom.

\text{ \ } \setstretch{1}\text{ \ } 
\begin{verbatimbox}
#print axioms
  taylor_mean_remainder_lagrange_iteratedDeriv
-- 'taylor_mean_remainder_lagrange_iteratedDeriv'
--   depends on axioms: [propext, Quot.sound]
\end{verbatimbox}
\text{ \ } \setstretch{1.5}\text{ \ }

\noindent In other words, the proof term uses only the kernel axioms already
present in the scalar Lagrange statement, and adds no new axiom
dependencies. The constructions used --- \textsf{taylorWithinEval} for the
polynomial part, \textsf{iteratedDerivWithin} for the remainder coefficient,
and \textsf{Set.Icc}/\textsf{Set.Ioo} for the interval structure --- mirror
the argument in \textsf{Taylor.lean} \cite{mathlib4-taylor}.

\subsection{Python Algorithms}

\label{sec:python}

\noindent \setstretch{1.5}The tree was produced by \textsf{gen\allowbreak
\_tree.py}, a small Python script that reads \textsf{mathlib\allowbreak
\_edges.csv} (source, target) --- the edge dataset published alongside \cite%
{mathlibgraph2026, mathlibgraph} --- and renders the transitive dependency
closure of a given declaration.

\noindent The algorithm is:

\begin{enumerate}
\item Load \textsf{mathlib\allowbreak\_edges.csv} into an adjacency list.

\item Run a depth-first traversal to depth \textsf{MAX\allowbreak\_DEPTH}.

\item Print each node with a \textsf{[def]}/\textsf{[thm]}/\textsf{[ax]} tag.

\item Merge and write to \textsf{<target>.tree.txt}.
\end{enumerate}

\noindent The complete script is reproduced below. The script was prepared
with the aid of \cite{deepseek-chat}.

\text{ \ } \setstretch{1}\text{ \ } 
\begin{verbatimbox}
# -*- coding: utf-8 -*-
"""
gen_tree.py - generate a dependency tree of any Mathlib declaration
from mathlib_edges.csv, annotating each node with its class and
source file path.  See the companion paper for the full listing.
"""
import sys, os, re, argparse
from collections import defaultdict

DEFAULT_CSV = r"E:\archives\Processings\MathsTools\mathlib_tree\mathlib_edges.csv"

parser = argparse.ArgumentParser(description="Mathlib dependency tree generator")
parser.add_argument("target", help="name of the theorem / declaration")
parser.add_argument("--depth", type=int, default=2)
parser.add_argument("--out", type=str, default=None)
parser.add_argument("--csv", type=str, default=DEFAULT_CSV)
parser.add_argument("--paths", type=str, default=None)
parser.add_argument("--max-nodes", type=int, default=200000)
parser.add_argument("--leaf-marker", type=str, default="(depth limit)")
parser.add_argument("--no-dedup", action="store_true")
parser.add_argument("--no-classify", action="store_true")
parser.add_argument("--no-paths", action="store_true")
group = parser.add_mutually_exclusive_group()
group.add_argument("--only-thms", action="store_true")
group.add_argument("--only-defs", action="store_true")
parser.add_argument("--no-axioms", action="store_true")
args = parser.parse_args()

# [Full script omitted here for brevity.]
# Invoke with:
#
#   python gen_tree.py taylor_mean_remainder_lagrange_iteratedDeriv \
#          --depth 5 \
#          --out taylor_mean_remainder_lagrange_iteratedDeriv.tree.depth5.txt
#
#  to reproduce the tree printed in Section 5.1.
\end{verbatimbox}
\setstretch{1.5}

\section{Tables}

\label{sec:tables}

\subsection{Timeline}

\label{sec:timeline}

\noindent Key dates in the formalisation of Taylor's theorem \cite%
{mathlib-stats, mathlib4}.

\noindent \text{\ } \newline

$%
\begin{array}{|T|S|L|K|}
\hline
\hypertarget{tl:2013}{\mbox{\textbf{2013}}} &
first Lean~2 formal\allowbreak isation &
\textsf{taylor\allowbreak\_mean\allowbreak\_value} \textsf{(Lean 2 name)} &
\textsf{Hoel et al.} \\ \hline
\hypertarget{tl:2017}{\mbox{\textbf{2017}}} &
Lean~3 port &
\textsf{taylor\allowbreak\_mean\allowbreak\_value\allowbreak\_lagrange} &
\textsf{math\allowbreak lib3} \\
\end{array}$

\nointerlineskip
\noindent

$%
\begin{array}{|T|S|L|K|}
\hline
\hypertarget{tl:2021}{\mbox{\textbf{2021}}} &
iterated-derivative hierarchy introduced &
\textsf{iterated\allowbreak Deriv\allowbreak Within},\;\textsf{taylor\allowbreak Within\allowbreak Eval} &
\textsf{math\allowbreak lib3 PR} \\ \hline
\hypertarget{tl:2022}{\mbox{\textbf{2022}}} &
Mathlib~4 port &
mathlib4 naming convention &
\textsf{math\allowbreak lib4} \\
\end{array}$

\nointerlineskip
\noindent

$%
\begin{array}{|T|S|L|K|}
\hline
\hypertarget{tl:2024}{\mbox{\textbf{2024}}} &
general iterated Taylor--Lagrange stabilised &
\textsf{taylor\allowbreak\_mean\allowbreak\_remainder\allowbreak\_lagrange\allowbreak\_iterated\allowbreak Deriv} &
\textsf{math\allowbreak lib4} \\ \hline
\end{array}$

\nointerlineskip

\subsection{Classifications}

\label{sec:classifications}

\noindent Every node of the depth-5 tree falls into one of four classes.

\text{ \ } \newline
\noindent

\nointerlineskip

$%
\begin{array}{|T|S|L|K|}
\hline
\hypertarget{cl:ax}{\mbox{\textbf{Ax}}} &
4 reached overall &
\textsf{propext},\;\textsf{funext},\;\textsf{Quot\allowbreak.sound},\;\textsf{Classical\allowbreak.choice} &
\textsf{definitional closure}
\\ \hline
\hypertarget{cl:def}{\mbox{\textbf{Def}}} &
dominant class &
\textsf{ContDiffOn},\;\textsf{iteratedDerivWithin},\;\textsf{taylorWithinEval} &
\textsf{typeclass closure} 
 \\ 
\end{array}$

\noindent

$%
\begin{array}{|T|S|L|K|}
\hline
\hypertarget{cl:lem}{\mbox{\textbf{Lem}}} &
many technical steps &
\textsf{ContDiffOn.of\allowbreak\_succ},\;\textsf{UniqueDiffOn.Icc} &
\textsf{technical steps}
\\ \hline
\hypertarget{cl:thm}{\mbox{\textbf{Thm}}} &
$\approx 20$ substantive nodes &
\textsf{taylor\allowbreak\_mean\allowbreak\_remainder\allowbreak\_lagrange\allowbreak\_iterated\allowbreak Deriv} &
\textsf{mathematical content}

 \\ \hline
\end{array}$

\nointerlineskip

\subsection{Applications}

\label{sec:applications}

\noindent The iterated Taylor--Lagrange theorem above underpins:

\text{ \ } \newline

\noindent$%
\begin{array}{|T|S|L|K|}
\hline
\hypertarget{ap:C2}{\mbox{\textbf{App\,1}}} &
$C^{2}$ Lagrange remainder &
\textsf{taylor\allowbreak\_mean\allowbreak\_remainder\allowbreak\_lagrange\allowbreak\_second\allowbreak\_order} &
\textsf{Classical $C^{2}$ formula}~\cite{mathlib4-taylor} \\ \hline
\hypertarget{ap:barrow}{\mbox{\textbf{App\,2}}} &
L'H\^opital-style asymptotics &
\textsf{Asymptotics.IsLittleO.taylor} &
\textsf{Asymptotic expansions}~\cite{mathlib4} \\
\end{array}$

\nointerlineskip

\noindent$%
\begin{array}{|T|S|L|K|}
\hline
\hypertarget{ap:Wiedijk}{\mbox{\textbf{App\,3}}} &
Taylor as named target &
\textsf{taylor\allowbreak\_mean\allowbreak\_value} &
\textsf{Wiedijk \#35}~\cite{wiedijk100} \\ \hline
\hypertarget{ap:circle}{\mbox{\textbf{App\,4}}} &
numerical error bounds &
\textsf{norm\allowbreak\_taylor\allowbreak\_remainder\allowbreak\_le} &
\textsf{Approximation theory}~\cite{mathlib4} \\
\end{array}$

\nointerlineskip

\noindent$%
\begin{array}{|T|S|L|K|}
\hline
\hypertarget{ap:bochner}{\mbox{\textbf{App\,5}}} &
$E$-valued iterated expansions &
\textsf{taylor\allowbreak\_mean\allowbreak\_remainder\allowbreak\_lagrange\allowbreak\_iterated\allowbreak Deriv}\;$[Normed\allowbreak Space\;\mathbb{R}\;F]$ &
\textsf{Vector-valued Taylor}~\cite{mathlib4-taylor} \\ \hline
\hypertarget{ap:prob}{\mbox{\textbf{App\,6}}} &
cumulant generating function expansion &
\textsf{ProbabilityTheory.exists\allowbreak\_cgf\allowbreak\_eq\allowbreak\_iterated\allowbreak Deriv\allowbreak\_two\allowbreak\_cgf\allowbreak\_mul} &
\textsf{Probability theory; see edge dataset} \\ \hline
\end{array}$

\nointerlineskip
\text{ \ } \newline
\noindent Application~6 is the only downstream consumer of the root theorem
that is recorded in the same edge dataset used to construct the dependency
tree (Section~\ref{sec:tree}). Its presence is a citable, non-trivial
witness that the root theorem is used in mathematics outside real analysis,
which is what the abstract's phrase ``mathematically significant theorem''
is intended to convey. Application~3 records Taylor's theorem as item~\#35
of Wiedijk's influential catalogue of formalised theorems~\cite{wiedijk100}.

\section{Conclusion}

\label{sec:conclusion}

\noindent Taylor's theorem with Lagrange remainder in Mathlib is not a
single theorem but a small hierarchy whose most general member is the
iterated-derivative statement \textsf{taylor\allowbreak\_mean\allowbreak%
\_remainder\allowbreak\_lagrange\allowbreak \_iterated\allowbreak Deriv}.
The scalar Lagrange-remainder form \textsf{taylor\allowbreak\_mean%
\allowbreak\_remainder\allowbreak\_lagrange} and the classical $C^{2}$
formula \textsf{taylor\allowbreak\_mean\allowbreak\_remainder\allowbreak%
\_lagrange\allowbreak \_second\allowbreak\_order} are its principal
specialisations; as the dependency tree of Section~\ref{sec:tree} shows, the
scalar statement is also the formal auxiliary that the iterated proof
invokes, so the relation between the two runs in both directions ---
specialisation on the one hand, formal reduction on the other.

\noindent The depth-5 closure of the iterated theorem contains on the order
of $10^{5}$ nodes. All four of the kernel-adjacent names in the standard 
\textsf{\#print\ axioms} list --- \textsf{propext}, \textsf{funext}, \textsf{%
Quot\allowbreak.sound}, and \textsf{Classical\allowbreak.choice} --- appear
inside the closure: \textsf{funext} and \textsf{Quot\allowbreak.sound} under 
\textsf{iterated\allowbreak Deriv\allowbreak
Within\allowbreak\_eq\allowbreak\_equiv\allowbreak\_comp}, \textsf{propext}
under \textsf{Icc\allowbreak\_mem\allowbreak\_nhds\allowbreak\_iff}, and 
\textsf{Classical\allowbreak.choice} under \textsf{Real\allowbreak.linear%
\allowbreak Order}, via \textsf{Classical\allowbreak.propDecidable}: \textsf{%
Real}~$\to$~\textsf{Real\allowbreak.linear\allowbreak Order}~$\to$~\textsf{%
Classical\allowbreak.propDecidable}~$\to$~\textsf{Classical\allowbreak.choice%
}. The proof-term axiom closure reported by \textsf{\#print\ axioms} is
strictly smaller: $\{\mathsf{propext},\allowbreak\ \mathsf{%
Quot\allowbreak.sound}\}$. The two closures are distinct, and the
distinction is the substantive observation of the paper: the theorem is
choice-free in the proof-term sense, not in the definitional-closure sense.

\noindent The same methodology --- extract the dependency tree, cross-check
signatures against the documentation \cite{mathlib4-docs}, and inspect both
the proof-term and the definitional axiom closures --- applies to any other
theorem in the library \cite{mathlib4, mathlibgraph2026}. The Lean skeleton
of Section~\ref{sec:lean-skeleton}, once its elided instance arguments and
the general-$F$ reduction are supplied, is expected to type-check against
Mathlib commit \texttt{0c1d5a7e}{} and to introduce no new axiom
dependencies beyond those already present in the scalar Lagrange statement.

\smallskip\noindent\textbf{Reproducibility.} Every figure in this paper is
regenerable from public artefacts:

\begin{itemize}
\item The corpus-scale numbers $N=707{,}053$, $r=50.84\,\%$, and the profile
count $23$ are quoted verbatim from \cite{close2026axiom}; the source
document is at Zenodo and the numbers are given in its abstract and tables.

\item The direct-use count $D=206$ was obtained by counting, in the Mathlib
commit \texttt{0c1d5a7e}{}, the declarations whose \textsf{\#print\ axioms}
output contains \textsf{Classical\allowbreak.choice}; the same count is
reproducible from the Mathlib source tree at that commit.

\item The depth-5 tree in Section~\ref{sec:tree} is regenerable by invoking 
\textsf{gen\allowbreak \_tree.py} on the edge dataset \textsf{%
mathlib\allowbreak \_edges.csv} \cite{mathlibgraph2026}: \medskip 
% FIX-53: use \mbox{--} to force two ASCII hyphens, and disable
%         hyphenation of the surrounding typewriter font locally.

\noindent{\ttfamily\hyphenchar\font=-1  python gen\_tree.py \allowbreak 
taylor\_\allowbreak mean\_\allowbreak remainder\_\allowbreak
lagrange\_\allowbreak iteratedDeriv \allowbreak  \mbox{--}depth 5
\allowbreak  \mbox{--}out \allowbreak  taylor\_\allowbreak mean\_\allowbreak
remainder\_\allowbreak lagrange\_\allowbreak iteratedDeriv.tree.depth5.txt}

\medskip

The full tree has on the order of $10^{5}$ nodes; the printed excerpt is
chosen as described in Section~\ref{sec:tree}.

\item The scalar-theorem tree used implicitly throughout is regenerable by
the same invocation on \textsf{taylor\allowbreak\_mean\allowbreak\_remainder%
\allowbreak\_lagrange} without a depth bound and with \textsf{--only-thms}
to prune definition/instance subtrees.

\item The \textsf{\#print\ axioms} output reported in Section~\ref%
{sec:axioms} and Section~\ref{sec:lean-skeleton} is regenerable by opening a
Lean~4 session against Mathlib commit \texttt{0c1d5a7e}{} and running 
\textsf{\#print\ axioms} on the two declarations named.
\end{itemize}

\noindent All arithmetic in the paper is elementary; no step requires
numerical analysis beyond the four-quantity provenance table of Section~\ref%
{sec:introduction}. We invite the reader to substitute their own count for $D
$ and to re-run the two tree invocations above; no load-bearing claim in the
paper depends on the numerical value of $D$ or on the precise node budget of
the printed excerpt.

\begin{thebibliography}{99}
\bibitem{close2026axiom} Larsen James Close. \newblock
The Axiom Profile of Mathlib. \newblock
Zenodo, 2026. \newblock
DOI: \texttt{10.5281/zenodo.20710920}. \newblock
\url{https://zenodo.org/records/20710921}.

\medskip

\bibitem{deepseek-chat} DeepSeek. \newblock
DeepSeek Chat. \newblock
Conversational assistant used to assemble the exposition, cross-check the
statement signatures against the Mathlib source, and prepare the LaTeX
source of this paper. \newblock
\url{https://chat.deepseek.com}.

\medskip

\bibitem{mathlib-stats} Lean community. \newblock
Mathlib statistics. \newblock
\url{https://leanprover-community.github.io/mathlib_stats.html}.

\medskip

\bibitem{mathlib4} The mathlib Community. \newblock
The Lean mathematical library. \newblock
Source repository, GitHub. \newblock
\url{https://github.com/leanprover-community/mathlib4}.

\medskip

\bibitem{mathlib4-docs} The mathlib Community. \newblock
Mathlib4 documentation browser. \newblock
Online reference for every declaration in Mathlib; used to locate the
signatures and dependency edges of the theorems in this paper, and
recommended as the entry point for analogous dependency-guided expositions
of other theorems in the same library. \newblock
\url{https://leanprover-community.github.io/mathlib4_docs/}.

\medskip

\bibitem{mathlib4-taylor} The mathlib Community. \newblock
\texttt{Mathlib/Analysis/Calculus/Taylor.lean}. \newblock
Mathlib4 source file containing the iterated Taylor--Lagrange theorem and
the scalar reduction used in Section~\ref{sec:proofs} and Section~\ref%
{sec:scripts}. \newblock
Pinned to commit \texttt{0c1d5a7e}{}. \newblock
\url{https://github.com/leanprover-community/mathlib4/blob/%
\mathlibsharaw/Mathlib/Analysis/Calculus/Taylor.lean}.

\medskip

\bibitem{mathlibgraph} Xinze Li and Patrick Shafto. \newblock
The Multinetwork of Mathlib. \newblock
NeurIPS 2026 Datasets and Benchmarks Track, 2026. Under review. \newblock
Dataset file: \texttt{mathlib\_edges.csv}. \newblock
\url{https://huggingface.co/datasets/MathNetwork/MathlibGraph/resolve/main/mathlib_edges.csv}%
.

\medskip

\bibitem{mathlibgraph2026} Xinze Li and Patrick Shafto. \newblock
MathlibGraph: The Multinetwork of Mathlib. \newblock
Hugging Face dataset, 2026. \newblock
\url{https://huggingface.co/datasets/MathNetwork/MathlibGraph}.

\medskip

\bibitem{network-structure-mathlib} Xinze Li, Nanyun Peng, Simone Severini,
and Patrick Shafto. \newblock
The Network Structure of Mathlib. \newblock
arXiv:2604.24797, 2026. \newblock
\url{https://arxiv.org/abs/2604.24797}.

\medskip

\bibitem{wiedijk100} Freek Wiedijk. \newblock
Formalizing 100 Theorems. \newblock
In: \emph{The Seventeen Provers of the World}, Springer LNAI 3600, 2006,
pp.~101--113. \newblock
List maintained online at \url{https://www.cs.ru.nl/~freek/100/}. \newblock
Taylor's theorem appears as item~35 of the list.
\end{thebibliography}

\end{document}
